% !TeX root = ../thesis.tex

%%%%%%%%%%%%%%%%%%%%%%%%%%%%%%%%%%%%%%%%%%
\thesischapter[ch:permapprox][0.7\textwidth]
  {``The only way of discovering the limits of the possible\\ is to venture a little way past them into the impossible.''}
  {--- Arthur C. Clarke, 1962}
  {Approximation of permutation p-values with permApprox}
\nocite{clarke1973profiles}
\chapterabstract{This chapter presents the second main methodological contribution of the thesis: permApprox, a framework for accurate and zero-free approximation of permutation p-values when only a limited number of permutations is feasible. It motivates the finite-resolution problem of empirical permutation p-values, reviews GPD-based tail modeling, and addresses a key limitation of existing approaches, namely zero p-values caused by bounded fitted tails. The proposed solution uses constrained GPD fitting to ensure valid extrapolation beyond the observed statistic. The chapter further describes threshold selection, the data-adaptive safety margin \(\varepsilon\), GPD parameter estimation, and simulation studies evaluating accuracy, robustness, and computational performance.}
\chaptercontents
%%%%%%%%%%%%%%%%%%%%%%%%%%%%%%%%%%%%%%%%%%

\section*{Formal chapter information}

\subsection*{Author contributions}

Stefanie Peschel developed the methodology presented in this chapter and the associated \texttt{permApprox} framework. She conducted all simulation and application studies, produced the R scripts, and created the \texttt{permApprox} R package, including its documentation. The accompanying paper was written by Stefanie Peschel, and all graphics contained therein were created by her. All work presented in this chapter was carried out under the supervision and guidance of Prof. Dr. Christian L. Müller. In addition, Prof. Dr. Anne-Laure Boulesteix, co-author of the accompanying paper, contributed valuable ideas in the early stages of the project, proofread the final versions, and provided suggestions for improving the content and design of the paper.

\subsection*{Code and data availability}

The code and analysis material used to generate the simulation results, tables, and figures presented in this chapter are included in the thesis repository archived on Zenodo \citep{Peschel2026thesisRepo}. Within this repository, the material belonging to the present chapter is provided in the folder \href{https://github.com/stefpeschel/phd_thesis/tree/main/ch06_permapprox}{\texttt{ch06\_permapprox}}. Some large intermediate files from simulation studies are not archived in all cases, but can be regenerated by rerunning the corresponding scripts provided in the repository.

The \texttt{permApprox} R package implementing the methods presented in this chapter is available at: \url{https://github.com/stefpeschel/permApprox}

\subsection*{Note on publication}

The work contained in this Chapter is also presented in the following preprint, where Stefanie Peschel is first author:

\begin{itemize}
\item Stefanie Peschel, Anne-Laure Boulesteix, Erika von Mutius, and Christian L. Müller. (2026). permApprox: A general framework for accurate permutation $p$-value approximation. \textit{arXiv}.
\end{itemize}

The preprint is made available under the arXiv Non-Exclusive Distribution License (version 1.0), which permits distribution of the work by arXiv while the authors retain copyright. The full license text is available at:\\
\url{https://arxiv.org/licenses/nonexclusive-distrib/1.0/license.html}

\subsection*{Note on additional results}

The results presented in this chapter are based on the published article described above. In contrast to the original publication, the chapter has been extended to provide additional analyses and methodological refinements that support the overall workflow presented here. 
In particular, selected simulation studies have been complemented and extended, and additional methodological components have been included where relevant. The real-data application discussed in the preprint is not reproduced here, as it is incorporated into the integrated analysis presented in Chapter~\ref{ch:application}.

%%%%%%%%%%%%%%%%%%%%%%%%%%%%%%%%%%%%%
\section*{Notation}
\label{sec:6perm:notation}
%%%%%%%%%%%%%%%%%%%%%%%%%%%%%%%%%%%%%

In this chapter, additional notation is needed, which has not been defined in the former sections. Table~\ref{tab:6perm:notation} gives an overview.

\begin{table}[H]
\centering
\caption{Notation used throughout this chapter.}
\label{tab:6perm:notation}
\begin{tabularx}{\textwidth}{lX}
\toprule
\textbf{Symbol} & \textbf{Meaning} \\
\midrule
\multicolumn{2}{l}{\textbf{General permutation test notation}} \\[2pt]
$n$ & Per-group sample size \\
$m$ & Number of hypothesis tests (features, genes, taxa) \\
$B$ & Number of permutations \\
$b \in \{1,\dots,B\}$ & Index for permutation replicates \\
$j \in \{1,\dots,m\}$ & Index for hypothesis tests \\
$T_{\text{obs}}^{(j)}$ & Observed test statistic for test $j$ \\
$T_b^{*(j)}$ & Permuted test statistic for test $j$ in permutation $b$ \\
$\mathcal{T}^{(j)} = \{T_b^{*(j)}\}_{b=1}^B$ & Permutation distribution for test $j$ \\
\midrule

\multicolumn{2}{l}{\textbf{$p$-values}} \\[2pt]
$p_{\mathrm{emp}}$ & Empirical permutation $p$-value \\
$p_{\mathrm{GPD}}$ & Unconstrained GPD-based $p$-value \\
$p_{\mathrm{cGPD}}$ & Constrained GPD-based $p$-value \\
$p$ & Hybrid $p$-value combining empirical and GPD-based $p$-values \\
\midrule

\multicolumn{2}{l}{\textbf{GPD tail modeling}} \\[2pt]
$p_{\mathrm{thr}}$ & GPD tail modeling threshold (on $p_{\mathrm{emp}}$) \\
$u$ & Threshold defining the GPD tail region \\
$k$ & Number of exceedances above $u$\\
$Y_{\text{obs}}$ & Observed excess above $u$ \\
$\mathcal{Y} = \{Y_b^*\}$ & Permutation excesses above $u$ \\
$\sigma, \xi$ & Scale and shape parameters of the GPD \\
$s = -\sigma/\xi$ & Upper support boundary when $\xi < 0$ \\
$\hat{\sigma}_{(U)}, \hat{\xi}_{(U)}$ & Unconstrained GPD parameter estimates \\
$\hat{\sigma}_{(C)}, \hat{\xi}_{(C)}$ & Constrained GPD parameter estimates \\
$\hat{s}_{(U)}, \hat{s}_{(C)}$ & Unconstrained / constrained support boundaries \\
$\bar{F}_{\mathrm{GPD}}(\cdot) = 1 - F_{\mathrm{GPD}}(\cdot)$ & Upper GPD tail probability\\
\midrule

\multicolumn{2}{l}{\textbf{Data-adaptive $\varepsilon$ rule (SLLS)}} \\[2pt]
$\varepsilon$ & Safety margin ensuring GPD support at $Y_{\text{obs}}$
\\
$\kappa_{\mathrm{factor}}$ & Tuning constant controlling curvature in the log-saturation component\\
$\tau$ & Tuning constant controlling plateau height \\
$\rho_{\mathrm{lift}}$ & Weighting parameter in lift component \\
\bottomrule
\end{tabularx}
\end{table}

%%%%%%%%%%%%%%%%%%%%%%%%%%%%%%%%%%%%%
\section{Introduction}
\label{sec:6perm:introduction}
%%%%%%%%%%%%%%%%%%%%%%%%%%%%%%%%%%%%%

Permutation tests provide a flexible and widely applicable framework for statistical inference when analytic reference distributions are unavailable or unreliable. 
Rather than relying on parametric assumptions, permutation procedures exploit exchangeability under the null hypothesis: by repeatedly reshuffling labels and recalculating a test statistic, an empirical approximation to its null distribution is obtained \citep{fisher1935design,pitman1937significance,good2005permutation}. 
This conceptual simplicity, together with minimal distributional requirements, has made permutation testing a central tool in modern high-dimensional data analysis.

In microbiome research, permutation-based inference plays an especially important role. 
Many statistics arising at different inferential targets, such as differential abundance testing, community-level comparisons, and network-based comparisons, lack tractable parametric null distributions. 
Moreover, complex data characteristics such as sparsity, compositional constraints, and heterogeneity often make classical model-based approximations unreliable. 
Permutation procedures therefore serve as a unifying inferential mechanism across feature-level, community-level, and structural analyses.

Despite their theoretical validity, permutation tests face practical limitations when accurate estimation of small $p$-values is required. Under a finite number $B$ of permutations, the empirical null distribution is discrete and the attainable $p$-values lie on a grid of resolution $1/(B+1)$. If the observed test statistic exceeds all permuted values, the empirical $p$-value becomes zero. 
Such zero $p$-values are formally invalid \citep{phipson2010permutation} and remain zero after multiplicity adjustment, potentially distorting downstream inference.

A commonly recommended remedy to avoid zero $p$-values is to include the observed test statistic in the permutation set \citep{north2002note,phipson2010permutation}, which guarantees strictly positive $p$-values. However, this approach raises a new problem: the smallest attainable $p$-value remains bounded below by $1/(B+1)$, which may be insufficient in large-scale multiple testing settings. Increasing $B$ improves resolution but can become computationally prohibitive, particularly when test statistics are costly to evaluate or when thousands of hypotheses are assessed simultaneously \citep{john2022efficient}. 
This limitation is especially pronounced in comparative network analysis as discussed in Chapter~\ref{ch:netcomi}, where constructing a single microbial network already involves multi-step estimation, sparsification, and summary computation, so that repeating this workflow for a sufficiently large number of permutations is often computationally infeasible.

% Copied from the paper:
To overcome the limitations of finite permutation resolution, \citet{knijnenburg2009fewer} introduced a parametric tail approximation based on the Generalized Pareto Distribution (GPD), which enables accurate extrapolation of small $p$-values from a limited number of permutations. Their method fits a GPD to the largest subset of permuted test statistics and uses the resulting model to extrapolate beyond the observed permutation range, thereby reducing the number of permutations needed while retaining accuracy in the small-$p$ regime. Building on this idea, \citet{winkler2016faster} conducted a systematic comparison of several strategies for accelerating permutation inference in the general linear model, including (i) early stopping with negative binomial modeling, (ii) Gamma distribution fitting to the full permutation distribution, and (iii) GPD fitting restricted to the tail. They recommended the GPD-based tail approximation as a general-purpose choice due to its flexibility and weak distributional assumptions. The authors also emphasized that including or excluding the observed test statistic in the permutation set can systematically shift the resulting $p$-values, leading either to conservativeness or to ``invalid'' $p$-values.

However, an important limitation remains insufficiently addressed. When the estimated GPD shape parameter $\xi$ is negative, the fitted distribution has a finite upper endpoint. If the observed test statistic lies beyond this estimated support boundary, the extrapolated $p$-value again collapses to zero. Ad hoc remedies, such as monotone transformations of the test statistic to enforce positive shape estimates \citep{knijnenburg2009fewer}, alter the scale of the statistic and may affect the results. Prior work has not provided a principled and general solution that guarantees strictly positive tail-approximated permutation $p$-values while preserving the original scale of the test statistic.

The methodological development presented in this chapter addresses precisely this gap. Building on the GPD-based $p$-value approximation approach proposed by \citet{knijnenburg2009fewer}, we introduce a \emph{support-constrained} parameter estimation strategy that explicitly enforces compatibility between the fitted tail model and the observed test statistic. By constraining the estimated support to extend beyond the observed value, the procedure guarantees strictly positive $p$-values while preserving the statistical structure of the original test statistic.

Beyond the core constraint mechanism, this chapter develops a comprehensive workflow for GPD-based permutation inference. Key practical decisions, including the choice of GPD parameter estimation method, selection of the tail threshold, and integration of hybrid empirical-parametric $p$-values into downstream inference procedures, are formalized within a unified framework. The workflow is implemented in the open-source \texttt{R} package \texttt{permApprox}, which provides modular implementations of the alternative approximation strategies discussed in this chapter and thereby facilitates systematic comparative evaluation.

The performance and robustness of the proposed methodology are evaluated in extensive simulation studies based on classical two-sample tests, including the Student’s $t$-test and the Wilcoxon rank-sum test, covering both parametric and nonparametric settings. The results demonstrate that the support-constrained GPD approach achieves accurate tail probability estimation, eliminates zero $p$-values, and yields stable behavior even under limited permutation budgets.

Conceptually, this development extends beyond any single application domain. 
While permutation testing arises naturally in comparative network analysis (Chapter~\ref{ch:netcomi}), its \text{limitations} affect comparative inference at all inferential targets covered in Chapters~\ref{ch:feature} to \ref{ch:netcomi}. The methodology developed here, therefore, strengthens the inferential foundation of permutation-based testing throughout the thesis.

%%%%%%%%%%%%%%%%%%%%%%%%%%%%%%%%%%%%%
\section{Permutation testing under finite resolution}
\label{sec:6perm:finite_resolution}
%%%%%%%%%%%%%%%%%%%%%%%%%%%%%%%%%%%%%

Permutation tests provide a flexible nonparametric framework for hypothesis testing when the sampling distribution of a test statistic is unknown or difficult to derive analytically \citep{good2005permutation}. 
The basic idea dates back to the work of \citet{fisher1935design}, who proposed randomization-based significance tests for experimental data. 
Subsequent developments by \citet{pitman1937significance} and others established the theoretical foundations of permutation inference. 
A detailed historical overview of permutation testing and its practical implementation is provided by \citet{phipson2010permutation}.

The key principle underlying permutation inference is exchangeability. Under the null hypothesis, the joint distribution of the observations is invariant to permutations of the labels that define the experimental groups. Consequently, recomputing the test statistic after permuting these labels generates realizations from the null distribution of the statistic \citep{good2005permutation}.

Let $T^{(j)}$ denote a test statistic associated with hypothesis test $j \in \{1,\dots,m\}$ and let $T_{\mathrm{obs}}^{(j)}$ denote the observed value computed from the original data. 
Throughout this chapter we consider right-tailed tests, where large values of the statistic provide evidence against the null hypothesis. 
For left-tailed or two-sided tests, the statistics can be transformed (for example by sign inversion or by taking absolute values) so that inference is always conducted on the upper tail.

In principle, permutation testing is based on the \emph{exact permutation distribution}, obtained by evaluating the test statistic for all possible permutations of the data labels \citep{phipson2010permutation}. 
Let $\mathcal{P}$ denote the set of all such permutations and let $|\mathcal{P}|$ denote its cardinality. For each permutation $b \in \{1,\dots,|\mathcal{P}|\}$, let $T_b^{*(j)}$ denote the value of the test statistic computed after applying the $b$-th permutation. 
The exact permutation distribution is then given by the set
\[
\{T_b^{*(j)}\}_{b=1}^{|\mathcal{P}|}.
\]
The exact permutation $p$-value is defined as the proportion of permutations for which the test statistic is at least as extreme as the observed value:
\begin{equation}
p^{(j)} =
\frac{\sum_{b=1}^{|\mathcal{P}|}
\mathbb{I}\!\left(T_b^{*(j)} \ge T_{\mathrm{obs}}^{(j)}\right)}
{|\mathcal{P}|}.
\end{equation}

As the number of possible permutations grows rapidly with sample size, using the complete set of permutations is usually not feasible in practice. For example, in a two-sample test with group sizes $n_1$ and $n_2$, the number of distinct permutations of the group labels equals
$|\mathcal{P}| = \binom{n_1+n_2}{n_1}$.
Even for moderate sample sizes such as $n_1=n_2=20$, the total number of permutations already exceeds $10^{11}$.

Instead, practical permutation tests evaluate the test statistic only for a randomly selected subset of $B$ permutations drawn from $\mathcal{P}$. 
Denoting the resulting permutation statistics by $\{T_b^{*(j)}\}_{b=1}^{B}$, the sampled permutation distribution is given by
\begin{equation}
\mathcal{T}^{(j)} = \{T_b^{*(j)}\}_{b=1}^{B}.
\end{equation}
This sampled distribution provides an approximation to the exact permutation distribution and forms the basis for the empirical permutation $p$-value used in practice \citep{phipson2010permutation}.

The strength of evidence against the null hypothesis is quantified through the empirical permutation $p$-value
\begin{equation}
\label{eq:6perm:empirical_pvalue}
p_{\mathrm{emp}}^{(j)} =
\frac{1 + \sum_{b=1}^{B}\mathbb{I}\!\left(T_b^{*(j)} \ge T_{\mathrm{obs}}^{(j)}\right)}{1 + B},
\end{equation}
where $\mathbb{I}(\cdot)$ denotes the indicator function. 
The addition of 1 to both the numerator and denominator ensures that the resulting $p$-value is strictly positive even when the observed statistic exceeds all sampled permutation values. 
This adjustment has been strongly recommended by \citet{phipson2010permutation}, since a permutation $p$-value cannot be exactly zero in principle: the exact permutation distribution necessarily includes the observed data configuration itself.

Although the correction in Equation~\eqref{eq:6perm:empirical_pvalue} prevents zero $p$-values, it does not eliminate a fundamental limitation of permutation inference in practice. 
When the permutation distribution is approximated using a finite number of sampled permutations, the attainable $p$-values lie on a discrete grid. 
Specifically, for $B$ sampled permutations the empirical $p$-values are restricted to multiples of $1/(B+1)$, so that the smallest attainable value equals $1/(B+1)$. 
Consequently, the resolution with which small tail probabilities can be represented is determined entirely by the number of permutations performed.

This finite-resolution property becomes problematic when the permutation budget is limited. 
For example, with $B=1000$ permutations the smallest attainable empirical $p$-value is $1/(B+1) \approx 10^{-3}$. Even if the exact permutation $p$-value were substantially smaller, the empirical approximation cannot represent values below this threshold. 
In large-scale multiple testing problems, however, substantially smaller $p$-values may be required for findings to remain significant after adjusting for multiplicity. 
Achieving such resolution would require a much larger number of permutations. 
In many practical applications, however, performing a sufficiently large number of permutations is computationally infeasible, particularly when the test statistic itself is expensive to evaluate or when thousands of hypotheses must be assessed simultaneously.

% Refer to the network comparison chapter here again?

%%%%%%%%%%%%%%%%%%%%%%%%%%%%%%%%%%%%%
\section{GPD-based \texorpdfstring{$p$}{p}-value approximation}
\label{sec:6perm:gpd_tail_modeling}
%%%%%%%%%%%%%%%%%%%%%%%%%%%%%%%%%%%%%

The finite-resolution limitation discussed in the previous section motivates the use of tail modeling techniques that allow small permutation $p$-values to be approximated beyond the empirical permutation grid. Rather than increasing the number of sampled permutations $B$, these approaches exploit the statistical structure of the extreme tail of the permutation distribution to approximate tail probabilities corresponding to rare events.

A principled framework for modeling extreme values is provided by extreme value theory. In particular, the peaks-over-threshold (POT) approach models the extreme tail of a distribution by focusing only on observations exceeding a sufficiently high threshold $u$. Instead of modeling the entire distribution, the POT framework considers the excesses above this threshold and approximates their distribution using a Generalized Pareto Distribution (GPD) \citep{pickands1975statistical,coles2001introduction}. This provides a flexible way to describe the tail behavior of a distribution without imposing strong assumptions on its overall shape.

As defined in the previous section, let $T_{\mathrm{obs}}^{(j)}$ and $\{T_b^{*(j)}\}_{b=1}^{B}$ denote the observed and permuted test statistics corresponding to hypothesis $j$. To model the tail of this permutation distribution, we select a threshold $u^{(j)}$ that defines the beginning of the tail region. The selection of an appropriate threshold is a central aspect of peaks-over-threshold modeling and will be discussed in detail in Section~\ref{sec:6perm:threshold}.

% Copied from paper:		
We refer to permutation values satisfying $T_b^{*(j)} > u^{(j)}$ as \emph{exceedances}. Conditional on exceeding the threshold, we define the associated \emph{excesses} as the amount by which the test statistic exceeds the threshold,
\begin{equation}
Y_b^{*(j)} = T_b^{*(j)} - u^{(j)} \;\big|\; T_b^{*(j)} > u^{(j)}, \quad b = 1, \dots, B.
\end{equation}
The collection of permutation excesses is denoted by
\begin{equation}
\mathcal{Y}^{(j)} = \{ Y_b^{*(j)} \}.
\end{equation}
The observed excess is defined analogously as
\begin{equation}
Y_{\mathrm{obs}}^{(j)} = T_{\mathrm{obs}}^{(j)} - u^{(j)} \;\big|\; T_{\mathrm{obs}}^{(j)} > u^{(j)}.
\end{equation}

According to the peaks-over-threshold framework of extreme value theory, the conditional distribution of exceedances above a sufficiently high threshold converges to a Generalized Pareto Distribution as the threshold increases \citep{pickands1975statistical,park2016estimating}. This result provides the theoretical basis for modeling the tail behavior of the permutation distribution using the excesses $\mathcal{Y}^{(j)}$.

The GPD is parameterized by a scale parameter $\sigma > 0$ and a shape parameter $\xi \in \mathbb{R}$, with cumulative distribution function
\begin{equation}
	\label{eq:6perm:GPD_CDF}
	F(y; \sigma, \xi) =
	\begin{cases}
		1 - \left(1 + \frac{\xi y}{\sigma} \right)^{-1/\xi}, & \text{if } \xi \neq 0, \\[6pt]
		1 - \exp\left(-\frac{y}{\sigma}\right), & \text{if } \xi = 0,
	\end{cases}
	\qquad y \geq 0,\quad 1+\frac{\xi y}{\sigma}>0.
\end{equation}
For $\xi<0$, the distribution has finite upper endpoint $-\sigma/\xi$, such that $F(y;\sigma,\xi)=1$ for $y\geq -\sigma/\xi$. In practice, the parameters $\sigma$ and $\xi$ must be estimated from the permutation excesses $\mathcal{Y}^{(j)}$. 
Different estimation approaches can be used for this purpose and will be discussed in detail in Section~\ref{sec:6perm:param_est}.

The shape parameter $\xi$ determines the tail behavior and support of the distribution \citep{hosking1987parameter}:
\begin{itemize}
    \item For $\xi>0$, the distribution has a heavy tail with unbounded support. 
    \item When $\xi=0$, the GPD reduces to the exponential distribution with mean $\sigma$.
    \item For $\xi<0$, the distribution has a finite upper support boundary at $s=-\sigma/\xi$, so that $F(y)=1$ for all $y \ge s$.
\end{itemize}
The bounded-support case is particularly relevant for permutation distributions, as it implies that the fitted tail distribution assigns zero probability beyond the boundary.

Note that the GPD appears in multiple parametrizations across the literature. We adopt the standard two-parameter convention from extreme value theory with scale $\sigma>0$ and shape $\xi\in\mathbb{R}$, which implies bounded support when $\xi<0$ \citep{coles2001introduction}. Alternative parametrizations, such as $(a,k)=(\sigma,-\xi)$, are also common, e.g., \citep{castillo2015likelihood}. The two formulations are directly related through this transformation.

Let $k^{(j)} = |\mathcal{Y}^{(j)}|$ denote the number of permutation statistics exceeding the threshold $u^{(j)}$. The empirical probability of exceeding the threshold is then given by
\[
p_{\mathrm{exc}}^{(j)} = \frac{k^{(j)}}{B}.
\]
Combining this exceedance probability with the tail probability of the fitted GPD yields an approximation of the permutation $p$-value,
\begin{equation}
\label{eq:6perm:GPD_pvalue}
p_{\mathrm{GPD}}^{(j)} =
p_{\mathrm{exc}}^{(j)}
\,
\bar{F}_{\mathrm{GPD}}\!\left(
Y_{\mathrm{obs}}^{(j)};
\hat{\sigma}^{(j)},\hat{\xi}^{(j)}
\right),
\end{equation}
where $\bar{F}_{\mathrm{GPD}} = 1 - F_{\mathrm{GPD}}$ denotes the upper tail probability of the fitted GPD.

GPD-based tail modeling for permutation testing was first proposed by \citet{knijnenburg2009fewer} as a computational strategy for approximating small permutation $p$-values using only a limited number of permutations. They demonstrated that this approach yields accurate $p$-values, even when the number of permutations is drastically reduced. Later work has confirmed this finding \citep{winkler2016faster}. However, as discussed in the following section, standard GPD fitting procedures may produce problematic tail approximations in certain situations, particularly when the estimated shape parameter is negative and the resulting support boundary lies below the evaluation point.


%%%%%%%%%%%%%%%%%%%%%%%%%%%%%%%%%%%%%
\section{Proposed constrained GPD fitting}
\label{sec:6perm:constraint}
%%%%%%%%%%%%%%%%%%%%%%%%%%%%%%%%%%%%%

The GPD-based tail approximation described in the previous section provides a principled way to approximate small permutation $p$-values beyond the empirical permutation grid. However, a practical difficulty arises when the fitted GPD has bounded support. As discussed above, this situation occurs when the estimated shape parameter is negative ($\hat{\xi}<0$), implying that the distribution has a finite upper endpoint. In this case, the tail approximation may produce a $p$-value equal to zero if the observed test statistic lies beyond this boundary. Figure~\ref{fig:6perm:GPD_standard_constrained} illustrates this issue together with the proposed remedy. For clarity, we omit the test index $j$ throughout this section and write all quantities for a fixed test.

To formalize this problem, let $(\hat{\sigma}, \hat{\xi})$ denote the parameter estimates obtained from fitting a GPD to the permutation excesses, as proposed by \citet{knijnenburg2009fewer}. When $\hat{\xi} < 0$, the fitted distribution has a finite upper support boundary given by
\begin{equation}
  \hat{s} = -\frac{\hat{\sigma}}{\hat{\xi}}.
\end{equation}

If the observed excess $Y_{\mathrm{obs}}$ lies at or beyond this boundary, the fitted tail distribution assigns zero probability mass to the evaluation point, resulting in a GPD-based $p$-value of zero. This situation is illustrated in the upper right panel of Figure~\ref{fig:6perm:GPD_standard_constrained}.

\begin{figure}[!ht]
  \centering
  \includegraphics[width=0.95\textwidth]{figures/permApprox/GPD_standard_constrained.png}
  \caption{
  \textbf{GPD-based $p$-value approximation.}
  \textbf{Left panel:} Distribution of the permuted test statistic $T^*$ with threshold $u$ defining the permutation excesses (orange), and observed test statistic $T_{\mathrm{obs}}$.
  \textbf{Upper right:} Standard (unconstrained) GPD tail approximation when the estimated shape parameter is negative ($\hat{\xi}<0$), resulting in a finite upper support boundary $\hat{s}$. If the observed excess $Y_{\mathrm{obs}}$ lies beyond this boundary, the tail approximation assigns zero probability mass, yielding a zero GPD-based $p$-value.
  \textbf{Lower right:} Proposed constrained GPD fit, enforcing that the upper support boundary $\hat{s}_{(C)}$ lies strictly above the evaluation point, with offset $\varepsilon$. This guarantees a strictly positive tail probability at $Y_{\mathrm{obs}}$.
  This figure is adapted from \citep{peschel2026permapprox}.
  }
  \label{fig:6perm:GPD_standard_constrained}
\end{figure}

Zero-valued $p$-values produced by the GPD tail approximation are undesirable for the same reasons discussed earlier for empirical permutation $p$-values.
As emphasized by \citet{phipson2010permutation}, permutation $p$-values should never be exactly zero, since the exact permutation distribution necessarily includes the observed data configuration. This principle should equally apply to approximations of permutation $p$-values obtained through tail modeling.

In practice, the bounded-support case with $\hat{\xi}<0$ occurs frequently and can therefore lead to zero-valued tail approximations. Empirical examples illustrating this behavior are provided in the single-cell RNA-seq analysis of \citet{peschel2026permapprox} (not shown in this thesis) and in the comparative analyses presented in Chapter~\ref{ch:application} of this thesis. These examples demonstrate that negative shape estimates and the resulting zero $p$-values are not merely theoretical edge cases but arise regularly in realistic data settings.

To address this issue, we propose to impose a support constraint during GPD parameter estimation. Specifically, the fitted distribution is required to have an upper support boundary that lies strictly above the evaluation point. 
Formally, we enforce the constraint
\begin{equation}
\fbox{$\displaystyle
\hat{s}_{(C)} > Y_{\mathrm{obs}} + \varepsilon
$}
\end{equation}
where $\varepsilon > 0$ denotes a small safety margin and the subscript $(C)$ refers to the constrained fit.

Under this constraint, the parameter estimates $(\hat{\sigma}_{(C)},\hat{\xi}_{(C)})$ are chosen such that the resulting support boundary
\[
\hat{s}_{(C)} = -\frac{\hat{\sigma}_{(C)}}{\hat{\xi}_{(C)}}
\]
lies strictly above the evaluation point. 
This guarantees that the fitted tail distribution assigns positive probability mass to $Y_{\mathrm{obs}}$. 
Consequently, the resulting tail probability remains strictly positive, as illustrated in the lower right panel of Figure~\ref{fig:6perm:GPD_standard_constrained}.

The constrained estimates are subsequently substituted into the GPD tail approximation formula in Equation~\eqref{eq:6perm:GPD_pvalue}, yielding the constrained GPD-based $p$-value $p_{\mathrm{cGPD}}$.

Conceptually, this procedure can be interpreted as restricting the feasible parameter space of the GPD to parameter combinations for which the fitted support includes the evaluation point. 
The constraint therefore prevents degenerate tail approximations while preserving the general flexibility of the GPD model. 
A key remaining component of this approach is the choice of the safety margin $\varepsilon$, which determines how far the support boundary must lie above the observed excess and thus influences the resulting tail probabilities. The construction of this margin is discussed in the following section.

%%%%%%%%%%%%%%%%%%%%%%%%%%%%%%%%%%%%%
\section{The \texttt{permApprox} workflow}
\label{sec:6perm:workflow}
%%%%%%%%%%%%%%%%%%%%%%%%%%%%%%%%%%%%%

%====================================
\subsection{Overview and scope of the workflow}
\label{sec:6perm:workflow:overview}
%====================================

Beyond the support-constrained GPD fitting approach introduced in the previous section, accurate permutation-based inference requires a sequence of additional methodological decisions that extend beyond the tail model itself. In particular, one must determine (i) which empirical permutation $p$-values should be refined by tail approximation, (ii) how the tail region of the permutation distribution is defined, (iii) which methods are suitable for GPD parameter estimation, and (iv) how the support constraint is incorporated into the estimation procedure. 

To address these interdependent decisions in a principled and reproducible manner, we propose the \texttt{permApprox} workflow, which is implemented in an R package of the same name. The workflow integrates constrained GPD fitting with data-driven choices for all preceding steps, resulting in a coherent pipeline for permutation-based $p$-value approximation.

Figure~\ref{fig:6perm:workflow} provides a schematic overview of the workflow. Each step is discussed in detail in the subsequent sections of this chapter. For every stage, we consider a range of candidate methods, evaluate their performance in dedicated simulation studies, and derive a recommended default configuration. These default choices are summarized in Section~\ref{sec:6perm:workflow:default_config} and highlighted in Figure~\ref{fig:6perm:workflow}. All alternative methods considered throughout this chapter are available in the \texttt{permApprox} package, allowing users to adapt the workflow to specific applications.

\subsubsection*{Input}

The workflow operates on a vector of observed test statistics $T_{\mathrm{obs}} = (T_{\mathrm{obs}}^{(1)}, \dots, T_{\mathrm{obs}}^{(m)})$ and a corresponding matrix of permuted test statistics $\mathcal{T} = \{T_b^{*(j)}\}$, where $b = 1, \dots, B$ indexes permutations and $j = 1, \dots, m$ indexes tests. In the special case $m = 1$, the input reduces to a single observed statistic and its permutation distribution. In the following, we focus on the multiple-testing setting.

\subsubsection*{Steps 1 and 2: screening for tail approximation}

The first decision concerns whether an empirical permutation $p$-value should be refined using parametric tail modeling. Empirical permutation $p$-values are exact under the null hypothesis, but their granularity is limited by the number of permutations $B$, which primarily affects inference in the extreme tail. Tail approximation is therefore applied only to sufficiently small empirical $p$-values. 

Restricting the approximation to the tail serves two purposes. First, it focuses modeling effort on the region where empirical estimates are most unreliable. Second, it avoids fitting a GPD to regions of the permutation distribution that are not well described by an extreme value model. For larger $p$-values, empirical estimates are sufficiently precise and are therefore retained.

Accordingly, the workflow applies a screening step (Step~1 in Figure~\ref{fig:6perm:workflow}) based on empirical permutation $p$-values. For each test $j$, the empirical value $p_{\mathrm{emp}}^{(j)}$ is computed as defined in Equation~\eqref{eq:6perm:empirical_pvalue}. Only tests with $p_{\mathrm{emp}}^{(j)} < p_{\mathrm{thr}}$ are selected for tail approximation (Step~2), where $p_{\mathrm{thr}}$ is a user-defined threshold. Following \citet{winkler2016faster}, we set $p_{\mathrm{thr}} = 2\alpha$ as a default choice.

\subsubsection*{Step 2a: threshold selection}

Since the GPD is intended to model only the extreme tail of a distribution, it is necessary to determine which permutation statistics belong to this region. This is achieved by selecting a threshold that separates the bulk of the permutation distribution from its tail. Observations exceeding this threshold are referred to as \emph{exceedances}, as introduced in Section~\ref{sec:6perm:gpd_tail_modeling}.

In Section~\ref{sec:6perm:threshold}, we investigate strategies for selecting an appropriate threshold for permutation distributions, including the choice of an initial number of exceedances $k_0$. The performance of these methods is evaluated in simulation studies with respect to both threshold recovery and downstream $p$-value approximation.

\subsubsection*{Step 2b: selection of the safety margin \texorpdfstring{$\varepsilon$}{epsilon}}

Given a threshold $u^{(j)}$ for each test $j$, the next step is to define the evaluation point at which the tail probability is approximated. In the proposed framework, this requires specifying a safety margin $\varepsilon^{(j)}$ and enforcing that the support of the fitted GPD includes the evaluation point $Y_{\mathrm{obs}}^{(j)} + \varepsilon^{(j)}$. 

In Section~\ref{sec:6perm:epsilon}, we develop a data-adaptive rule for selecting $\varepsilon^{(j)}$ based on the observed test statistics and their permutation distributions. This rule is designed to balance numerical stability and approximation accuracy across a wide range of settings.

In rare cases involving extremely large test statistics, numerical underflow may still occur, such that theoretically positive $p$-values are rounded to zero due to limited machine precision. To mitigate this issue, an optional refinement procedure for $\varepsilon$ is introduced in Section~\ref{sec:6perm:epsilon:refinement}.

\subsubsection*{Step 2c: GPD parameter estimation}

The final tail approximation step is the estimation of the GPD parameters from the exceedances while enforcing the support constraint. In Section~\ref{sec:6perm:param_est}, we introduce several candidate estimation methods and adapt them to the constrained setting.

Although all considered estimators are, in principle, compatible with the constrained framework, their performance differs substantially depending on sample size and the degree of misspecification induced by the constraint. To identify a suitable default method, we conduct simulation studies that evaluate both parameter recovery in well-specified settings and stability of tail estimates under constraint-induced misspecification.

For each test $j$, the output of this step consists of constrained parameter estimates $(\hat{\sigma}^{(j)}_{(C)}, \hat{\xi}^{(j)}_{(C)})$, obtained from a restricted parameter space such that
\begin{equation}
\hat{s}^{(j)}_{(C)} = -\hat{\sigma}^{(j)}_{(C)} / \hat{\xi}^{(j)}_{(C)} > Y_{\mathrm{obs}}^{(j)} + \varepsilon^{(j)}.
\end{equation}

\subsubsection*{Hybrid \texorpdfstring{$p$}{p}-values}

For tests selected for tail approximation, the constrained GPD fitting approach is used to compute refined tail probabilities $p_{\mathrm{cGPD}}^{(j)}$, which replace the corresponding empirical values. Tests that are not selected retain their empirical permutation $p$-values. This leads to a hybrid construction that combines empirical and parametric components.

Specifically, for each test $j$, the hybrid $p$-value is defined as
\begin{equation}
\label{eq:6perm:hybrid_pvalue}
p^{(j)} :=
\begin{cases}
p_{\mathrm{emp}}^{(j)}, & \text{if } p_{\mathrm{emp}}^{(j)} \ge p_{\mathrm{thr}}, \\[6pt]
p_{\mathrm{cGPD}}^{(j)}, & \text{if } p_{\mathrm{emp}}^{(j)} < p_{\mathrm{thr}}.
\end{cases}
\end{equation}

The resulting vector $\{p^{(j)}\}_{j=1}^m$ combines discrete empirical values in the bulk of the distribution with continuous tail approximations in the extreme region. These hybrid $p$-values form the primary output of the \texttt{permApprox} workflow and serve as the basis for subsequent inference, including multiple testing correction.

\begin{figure}[H]
  \centering
  \includegraphics[width=\textwidth]{figures/permApprox/workflow.pdf}
  \caption{\textbf{Graphical representation of the \texttt{permApprox} workflow.}\\
  \textbf{Input:} A vector with observed test statistics $T_{\mathrm{obs}}$ and a matrix with permuted test statistics $\mathcal{T} = {T_b^{*(j)}}$ for $B$ permutations and $m$ tests.
  \textbf{Step 1:} Empirical permutation $p$-values are computed and screened for small values relative to the chosen significance level $\alpha$ (i.e., $p_{\mathrm{emp}} < 2\alpha$), with selected tests flagged for GPD-based tail approximation.
  \textbf{Step 2:} For a selected test $j$, the tail approximation stage consists of three main steps:
  \textbf{Step 2a:} A threshold $u^{(j)}$ is chosen via the Anderson--Darling (AD) based test to define the tail region of the permutation distribution.
  \textbf{Step 2b:} The safety margin $\varepsilon^{(j)}$ is selected (with optional refinement in rare cases of machine-underflow). 
  \textbf{Step 2c:} A GPD is fitted to the permutation excesses under a support constraint ensuring valid evaluation beyond the observed excess.
  The fitted constrained GPD yields refined tail $p$-values $p_{\mathrm{cGPD}}^{(j)}$, which replace the corresponding empirical values.
  \textbf{Output:} Vector of hybrid $p$-values that are discrete in the bulk of the distribution and continuous in the extreme tail. The \textbf{red flags} mark our suggested default methods as explained in Section~\ref{sec:6perm:workflow:default_config}.
  This figure is adapted from \citep{peschel2026permapprox}.
  }
  \label{fig:6perm:workflow}
\end{figure}

%====================================
\subsection{Overview of simulation studies}
\label{sec:6perm:workflow:simulation_studies}
%====================================

To determine appropriate default settings for the \texttt{permApprox} workflow, a series of simulation studies was conducted, each addressing a specific component of the procedure. An overview of these studies and their interdependencies is provided in Figure~\ref{fig:6perm:sim_overview}, where the numbering reflects the order in which the studies were conducted.

While the workflow itself proceeds sequentially from threshold selection to parameter estimation and finally $p$-value approximation, the corresponding simulation studies were designed in a different order, reflecting both the relative importance of the individual components and their methodological dependencies.

In particular, the investigation began with GPD parameter estimation, as reliable estimation of the shape and scale parameters is fundamental for all subsequent steps. We first assessed the ability of six different estimators to recover the true GPD parameters in an unconstrained setting, followed by an evaluation of their robustness under support constraints, where parameter recovery is no longer meaningful (Sections~\ref{sec:6perm:param_est:sim_unconstrained} and \ref{sec:6perm:param_est:sim_constrained}). Based on these results, the likelihood moment estimator (LME) was selected as a suitable default and was used in all subsequent studies.

Building on this choice, we next investigated the influence of the initial number of exceedances $k_0$ used to start the threshold search (Section~\ref{sec:6perm:threshold:sim_exceed0}). This analysis was conducted using the classical failure-to-reject (FTR) method in combination with LME, and aimed to identify a stable and robust configuration for the permutation setting. Based on these results, a default choice of $k_0 = \max(0.25B, 250)$ was selected and used in all subsequent studies.

Subsequently, six different threshold selection methods were compared with respect to their ability to recover appropriate thresholds and to yield accurate $p$-value approximations (Sections~\ref{sec:6perm:threshold:sim_thresh_recovery} and \ref{sec:6perm:threshold:sim_pval_recovery}). These evaluations were carried out for a Gaussian $t$-test setting, allowing direct comparison to the corresponding $t$-test $p$-value as reference. Based on this comparison, the robust FTR method was selected as the default threshold detection approach and used in the final workflow.

Finally, the selected components were combined into the full \texttt{permApprox} workflow and compared to four competing approaches for $p$-value approximation (Section~\ref{sec:6perm:simulations}). This comparison was performed both for the $t$-test and for the Wilcoxon test applied to exponential data, thereby assessing the performance of the method under parametric as well as non-parametric settings.

\begin{figure}[H]
  \centering
  \includegraphics[width=\textwidth]{figures/permApprox/simulation_studies.pdf}
  \caption{\textbf{Simulation studies for selecting default \texttt{permApprox} settings.}\\
  Rectangular boxes represent components of the workflow, while rounded boxes indicate the objectives of the corresponding simulation studies. Dashed arrows denote the simulation studies, whereas solid arrows indicate the selected default configuration that is carried forward to subsequent steps. The numbering reflects the order in which the simulation studies were conducted, which differs from the sequential order of the workflow. The studies comprise evaluation of GPD parameter estimation methods (parameter recovery and robustness under constraint), selection of the initial number of exceedances $k_0$, comparison of threshold selection methods (threshold and $p$-value recovery), and final benchmarking of the complete workflow against competing approaches for $p$-value approximation.
  }
  \label{fig:6perm:sim_overview}
\end{figure}


%====================================
\subsection{Default \texttt{permApprox} configuration and implementation}
\label{sec:6perm:workflow:default_config}
%====================================

The final configuration of the \texttt{permApprox} workflow represents the outcome of several dedicated simulation studies, each addressing one specific component of the procedure. These studies were conducted to identify the most accurate, robust, and computationally efficient settings for practical use. Our proposed configuration is as follows:

\begin{itemize}

    \item \textbf{Initial number of exceedances:} The threshold search starts with the top 25\% of permutation values ($k_0 = 0.25B$), but includes at least 250 exceedances to ensure sufficient data for stable fitting. This configuration yielded the lowest bias and variance across simulated scenarios (Section~\ref{sec:6perm:threshold:sim_exceed0}).

    \item \textbf{GPD threshold detection:} The robFTR method in combination with the Anderson--Darling goodness-of-fit test is employed for threshold selection, as it achieved the best trade-off between robustness to outliers and computational efficiency. A detailed comparison of threshold detection strategies is provided in Section~\ref{sec:6perm:threshold}. The empirical $p$-value is used as fall-back, when the test is rejected for all candidate thresholds. 
      
    \item \textbf{Estimator:} We use the likelihood moment estimator (LME) as the default method for GPD parameter fitting. While any of the estimators evaluated in this thesis could, in principle, be applied, LME provides a favorable balance between estimation accuracy, robustness under support constraints, and computational efficiency in our simulation studies (Section~\ref{sec:6perm:param_est}).
    
    \item \textbf{Constraint:} The GPD is fitted under the boundary constraint described in Section~\ref{sec:6perm:epsilon}, using the standardized lifted log-saturation (SLLS) rule for determining $\varepsilon$. Throughout this thesis, the data-adaptive epsilon refinement procedure (Section~\ref{sec:6perm:epsilon:refinement}) is applied only in the real-world applications, but not in any of the simulation studies.

\end{itemize}

The proposed workflow is implemented in the open-source \texttt{R} package \texttt{permApprox}, which can be accessed at \url{https://github.com/stefpeschel/permApprox}.
The above configuration defines the package’s default settings, and was applied in the simulation study presented in Section~\ref{sec:6perm:simulations}, as well as the integrated analysis of the PASTURE cohort (Chapter~\ref{ch:application}. 
Additionally, the package provides all of the alternative $p$-value approximation, GPD estimation, and threshold detection methods evaluated in this work. This allows users to reproduce, compare, and extend our analyses.

%====================================
\subsection{Remarks on single and multiple testing}
\label{sec:6perm:workflow:multiple_testing}
% Section copied from the paper
%====================================

The \texttt{permApprox} framework can be applied both in single-test settings and in large-scale inference problems involving many parallel hypotheses. In the single-test setting, the output is a single hybrid $p$-value, whereas in the multiple testing setting it is a vector ${p^{(j)}}_{j=1}^m$ of hybrid $p$-values as defined in \eqref{eq:6perm:hybrid_pvalue}.

In multiple testing scenarios, the discreteness of empirical permutation $p$-values has motivated the development of several false discovery rate (FDR) procedures tailored specifically to discrete $p$-values \citep{gilbert2005modified,jiang2017discrete}. While classical procedures such as the Benjamini-Hochberg (BH) method  \citep{benjamini1995controlling} remain valid in this setting, they can exhibit conservative behavior because discrete null $p$-values tend to be stochastically larger than uniformly distributed ones \citep{benjamini2001control}.

The hybrid $p$-values produced by \texttt{permApprox} retain this discrete structure for moderate and large values, while they provide a smooth, continuous approximation in the extreme tail, where fine resolution is most critical. As a result, the hybrid $p$-values can often be treated analogously to continuous $p$-values in downstream inference, i.e., standard multiple testing procedures for continuous tests may be applied directly.

An important distinction between single and multiple testing settings arises in the selection of the cap $Z_{\mathrm{cap}}$ used in the data-adaptive $\varepsilon$ rule (Section~\ref{sec:6perm:epsilon:cap}). In multiple testing scenarios, the cap is defined by the maximum observed test statistic across all hypotheses, yielding a conservative and robust constraint. For single tests or small numbers of hypotheses, a permutation-based cap is used instead, providing greater stability in the presence of extreme observations.

%%%%%%%%%%%%%%%%%%%%%%%%%%%%%%%%%%%%%
\section{Threshold selection for permutation tails}
\label{sec:6perm:threshold}
%%%%%%%%%%%%%%%%%%%%%%%%%%%%%%%%%%%%%

Within the \texttt{permApprox} workflow introduced in Section~\ref{sec:6perm:workflow}, threshold selection constitutes the first step of the tail modeling stage (Step~2a). Its purpose is to identify the subset of permutation statistics that belong to the extreme tail and are therefore suitable for approximation by a generalized Pareto distribution (GPD).

The threshold choice plays a central role in extreme value modeling. If the threshold is chosen too low, the GPD is fitted to observations that do not follow the limiting tail behavior, leading to model misspecification and biased parameter estimates. Conversely, if the threshold is chosen too high, only few observations remain, resulting in high variability of the estimates. This bias-variance trade-off is particularly pronounced in permutation settings, where the number of available tail observations is inherently limited by the number of permutations. In this section, we review existing approaches to threshold selection and investigate additional methods that were developed in the course of this thesis.

In addition to evaluating the threshold selection methods themselves, the choice of the initial number of exceedances used to start the threshold search requires careful consideration in the permutation setting. In classical extreme value applications, it is common to select thresholds corresponding to approximately $5\%$ to $10\%$ of the sample, resulting in a moderate number of exceedances that balances bias and variance \citep{coles2001introduction}. In contrast, \citet{knijnenburg2009fewer} suggest using a fixed number of 250 exceedances in the context of permutation testing, which corresponds to $25\%$ of the data when $B=1000$ permutations are used. This discrepancy highlights that standard EVT heuristics may not directly translate to permutation-based settings. Therefore, in addition to comparing threshold selection methods, we systematically investigate the influence of the initial number of exceedances $k_0$ on the accuracy of resulting $p$-value approximations. Throughout this section, we use the term \emph{exceedances} to refer to observations above the threshold, and do not explicitly distinguish between exceedances and the corresponding excesses unless required for clarity.

%====================================
\subsection{Principles of threshold selection in POT modeling}
\label{sec:6perm:threshold:principles}
%====================================

In the peaks-over-threshold (POT) framework, the tail of a distribution is modeled by fitting a GPD to observations exceeding a \emph{sufficiently high} threshold. 
The theoretical justification for this approach follows from classical results in extreme value theory, which state that, for distributions lying in the domain of attraction of an extreme-value distribution, the conditional distribution of exceedances converges to a GPD as the threshold increases \citep{pickands1975statistical,park2016estimating}. 

In practice, however, the threshold cannot be chosen arbitrarily high. Selecting an appropriate threshold involves a fundamental bias-variance trade-off. If the threshold is chosen too low, the exceedances may still contain observations from the bulk of the distribution, violating the GPD assumption and introducing bias in the fitted model. Conversely, if the threshold is chosen too high, only a small number of exceedances remain, leading to high variance in the parameter estimates. Consequently, threshold selection aims to identify a region where the GPD approximation is sufficiently accurate while retaining enough exceedances for reliable parameter estimation \citep{knijnenburg2009fewer,scarrott2012review}.

An important structural property of the POT framework facilitates this search. Once the threshold is sufficiently high for the GPD approximation to hold, the excess distribution above any higher threshold also follows a GPD with the same shape parameter, albeit with fewer exceedances \citep{coles2001introduction,thompson2009automated}. This stability property can be expressed formally as follows. 
If exceedances above a threshold $u$ follow a GPD with parameters $(\sigma_u,\xi)$, then exceedances above any higher threshold $v>u$ also follow a GPD with the same shape parameter $\xi$ but with an adjusted scale parameter
\begin{equation}
\sigma_v = \sigma_u + \xi (v-u).
\end{equation}
Thus, while the number of exceedances decreases as the threshold increases, the shape parameter governing tail behavior remains invariant. 
This property provides the theoretical basis for diagnostic procedures that evaluate the adequacy of the GPD approximation across increasing candidate thresholds, which therefore typically examine a sequence of increasing thresholds and assess whether the GPD assumption appears plausible for the corresponding exceedances.

A variety of methods for threshold selection have been proposed in the extreme value literature \citep{scarrott2012review,langousis2016threshold}. Classical approaches rely on graphical diagnostics, such as mean residual life plots \citep{davison1990models} or parameter stability plots, which aim to visually identify regions where the GPD approximation becomes stable \citep{scarrott2012review}. These methods exploit the threshold stability property of the GPD and search for ranges of thresholds where diagnostic quantities exhibit approximately constant behavior. In addition to such graphical tools, automated procedures have been developed that aim to provide objective and reproducible threshold detection. For example, Goodness-of-Fit (GoF) statistics can be used to assess whether the exceedances above a candidate threshold are consistent with the GPD model \citep{scarrott2012review,langousis2016threshold}. Compared with graphical diagnostics, such automated procedures avoid subjective interpretation and are therefore particularly attractive when threshold selection must be applied repeatedly across many datasets.

%====================================
\subsection{Threshold selection in \texorpdfstring{$p$}{p}-value approximation}
\label{sec:6perm:threshold:role}
%====================================

While threshold selection has been extensively studied in extreme value theory, existing methods were primarily developed for classical applications such as environmental extremes or financial risk analysis. In the context of permutation-based $p$-value approximation, several additional practical requirements arise.

First, the threshold detection procedure must be applied repeatedly across a potentially large number of hypothesis tests. Consequently, the method must operate in a fully automated manner and avoid subjective decisions such as visual inspection of diagnostic plots. Second, the procedure must remain computationally efficient, as threshold selection and GPD fitting are performed for each tested hypothesis. Third, the method should be robust across a wide range of permutation distributions arising from different sample sizes, effect sizes, and numbers of permutations.

The original implementation of GPD-based permutation $p$-value approximation proposed by \citet{knijnenburg2009fewer} addressed the threshold problem using a simple heuristic procedure. In this approach, an initial tail region consisting of the largest 250 permutation values is defined and the threshold is subsequently increased in increments of ten observations. For each candidate threshold, a GPD is fitted and the adequacy of the model is evaluated using a GoF test. While this strategy provides a pragmatic and computationally efficient solution, it relies on fixed choices for the number of exceedances and threshold increments.

In this work, we therefore investigate a range of automated threshold selection strategies that satisfy the above design criteria. The considered methods include approaches motivated by existing diagnostics in extreme value analysis as well as additional procedures developed during the construction of the \texttt{permApprox} framework. Their performance is evaluated in dedicated simulation studies that assess both threshold selection behavior and the accuracy with which the resulting GPD approximations recover reference $p$-values across varying sample sizes, effect sizes, and numbers of permutations.

All considered methods operate on an increasing sequence of candidate thresholds in the upper tail of the permutation distribution. This sequence is initialized by specifying a starting number of exceedances $$k_0$$, which determines the lowest candidate threshold and hence the initial tail region considered for GPD fitting. For each subsequent candidate threshold, a GPD is fitted to the corresponding exceedances and evaluated using a method-specific selection criterion. The role of $$k_0$$ and its choice are examined separately in Section~\ref{sec:6perm:threshold:sim_exceed0}.

%====================================
\subsection{Existing threshold selection approaches}
\label{sec:6perm:threshold:existing}
%====================================

Among automated threshold selection procedures, GoF-based methods play a central role in extreme value analysis. These approaches evaluate whether the exceedances above a candidate threshold are consistent with the GPD model and select the smallest threshold for which the GoF test no longer indicates lack of fit.

A widely used heuristic of this type is the \emph{failure-to-reject} (FTR) rule proposed by \citet{choulakian2001goodnessoffit}. In this approach, candidate thresholds are evaluated sequentially along the upper tail of the sample. For each candidate threshold, a GPD is fitted to the corresponding exceedances and a GoF test is applied. The selected threshold is the smallest value for which the null hypothesis of the GoF test is not rejected. 
Formally, let $u_1 < u_2 < \dots < u_m$ denote an increasing sequence of candidate thresholds and let $p_i$ denote the $p$-value of the GoF test for the GPD fit at threshold $u_i$. 
The FTR rule selects the threshold
\begin{equation}
u = \min \{ u_i : p_i > \alpha \},
\end{equation}
where $\alpha$ denotes the significance level of the GoF test.
The underlying rationale is that once the threshold enters a region where the GPD approximation provides an adequate description of the tail, the GoF test should no longer detect significant deviations from the model.

GoF testing for the GPD is commonly performed using statistics such as the Anderson--Darling (AD) or Cramér--von Mises (CvM) tests \citep{choulakian2001goodnessoffit}. In particular, the AD statistic has been shown to have better sensitivity  against outliers in the tail than the CvM statistic, making it a popular choice in POT modeling.

The threshold search procedure proposed by \citet{knijnenburg2009fewer} for permutation $p$-value approximation can be interpreted as a practical implementation of the FTR principle. Their method evaluates a sequence of candidate thresholds along the permutation tail and accepts the first threshold for which the GPD fit is deemed adequate according to the AD test statistic.

While the FTR rule provides a simple and computationally efficient procedure, it may occasionally accept thresholds prematurely due to random fluctuations in the sequence of GoF $p$-values. To address this issue, more formal sequential testing procedures have been proposed.

One such approach is based on the \emph{ForwardStop} rule introduced by \citet{gsell2016sequential} for ordered hypothesis testing and later adapted to threshold selection in extreme value analysis by \citet{bader2018automated}. In this framework, the sequence of GoF tests performed at increasing candidate thresholds is viewed as a sequence of ordered hypotheses.

Let $u_1 < u_2 < \dots < u_m$ denote an increasing sequence of candidate thresholds and consider the hypotheses
\[
H_i: \text{The exceedances above threshold } u_i \text{ follow a GPD}, \qquad i = 1,\dots,m.
\]
For each threshold $u_i$, a GoF test is performed and the corresponding $p$-value is denoted by $p_i$.

Since the hypotheses are ordered, rejecting $H_i$ implies that all hypotheses corresponding to lower thresholds should also be rejected. This structure naturally leads to a sequential testing problem. The ForwardStop procedure transforms the sequence of $p$-values into cumulative statistics and determines the largest index $k$ for which the average transformed $p$-values remain below a predefined significance level $\alpha$.

Formally, the ForwardStop stopping index is defined by \citet{bader2018automated} as
\begin{equation}
\label{eq:6perm:thresh_fwdstop}
\hat{k}_{\mathrm{F}} = \max_{\substack{1 \le k \le m}}
\left\{
k : -\frac{1}{k} \sum_{i=1}^{k} \log(1-p_i) \le \alpha 
\right\}.
\end{equation}

Under suitable assumptions, this rule controls the false discovery rate for the ordered sequence of hypotheses \citep{gsell2016sequential}. In the context of threshold detection, selecting the stopping index $\hat{k}_{\mathrm{F}}$ corresponds to rejecting the hypotheses $H_1,\dots,H_{\hat{k}_{\mathrm{F}}}$ and accepting the first non-rejected threshold $u_{\hat{k}_{\mathrm{F}}+1}$ as the selected threshold. Intuitively, this procedure identifies the point along the sequence of candidate thresholds where the GoF tests begin to consistently support the GPD approximation.

To illustrate the behavior of these procedures, we consider a simple example based on a permutation test for the two-sample $t$-statistic. 
Two samples of size $n_1=n_2=250$ are generated from normal distributions with equal variance $\sigma=1$ and means $\mu_1=1$ and $\mu_2=0$, corresponding to a moderate effect size. A total of $B=1000$ permutations are used to obtain the permutation distribution of the test statistic.

Figure~\ref{fig:6perm:thresh_ftr_fwdstop} shows the permutation distribution together with the sequence of candidate thresholds evaluated by the threshold detection procedures. 
The candidate thresholds correspond to the ordered permutation statistics and are evaluated sequentially along the upper tail. 
To ensure that a sufficient number of exceedances is available for reliable GPD fitting, the search excludes the largest permutation values and requires at least 20 exceedances above the threshold (i.e., $\texttt{exceed\_min}=20$).

For each candidate threshold, a GPD is fitted to the corresponding exceedances and the Anderson--Darling GoF test is applied. Panel (b) in Figure~\ref{fig:6perm:thresh_ftr_fwdstop} displays the resulting sequence of AD $p$-values, while panel (c) shows the ForwardStop statistic (Equation~\eqref{eq:6perm:thresh_fwdstop}) computed from these $p$-values. The vertical lines in panel (b) indicate the thresholds selected by the FTR and ForwardStop rules. While the FTR rule selects the first candidate threshold for which the GoF test is not rejected, the ForwardStop procedure evaluates the entire sequence of tests before determining the stopping point.

In this work, we consider both the failure-to-reject (FTR) rule and the ForwardStop procedure as automated strategies for threshold detection. In Sections~\ref{sec:6perm:threshold:sim_thresh_recovery} and \ref{sec:6perm:threshold:sim_pval_recovery}, the performance of the methods is compared with that of additional threshold selection strategies proposed in this work. These strategies are evaluated using simulation experiments designed to assess the performance of the methods in threshold recovery as well as $p$-value recovery.

\begin{figure}[H]
\centering
\includegraphics[width=\textwidth]{figures/permApprox/thresh_ttest_ftr_fwd_stop_ADp_diag-1.pdf}
  \caption{
  \textbf{Threshold selection with FTR and ForwardStop.}
  Illustrative example based on two samples drawn from Gaussian distributions ($\mu_1=1$, $\mu_2=0$, $\sigma=1$) with per-group sample sizes $n_1=n_2=250$. The permutation test is based on the Student's $t$-statistic with $B=1000$ permutations. 
  \textbf{(a)} Permutation distribution of the test statistic. 
  \textbf{(b)} Anderson--Darling (AD) goodness-of-fit $p$-values obtained from fitting a GPD to the exceedances above each candidate threshold. The dashed horizontal line indicates the significance level $\alpha$ of the AD test, while the dotted vertical line marks the largest admissible threshold ensuring that at least 20 exceedances remain for fitting. The colored vertical lines indicate the thresholds selected by the FTR and ForwardStop methods. 
  \textbf{(c)} ForwardStop statistic evaluated across the sequence of candidate thresholds. The dashed horizontal line indicates the stopping criterion used by the method.
}
\label{fig:6perm:thresh_ftr_fwdstop}
\end{figure}

%====================================
\subsection{Additional approaches developed in this work}
\label{sec:6perm:threshold:proposed}
%====================================

While the existing methods described above provide automated strategies for threshold detection, their behavior may still be sensitive to random fluctuations in the sequence of GoF $p$-values. In particular, the failure-to-reject (FTR) rule may accept a threshold prematurely if a single GoF test fails to reject the null hypothesis by chance. To reduce the impact of such isolated fluctuations, we consider several additional threshold selection strategies that were explored during the development of the \texttt{permApprox} framework.

%------------------------------------
\subsubsection*{Robust failure-to-reject (robFTR)}
%------------------------------------

The first proposed approach is a conservative variant of the FTR rule. It requires several consecutive candidate thresholds to pass the GoF test before a threshold is accepted. The motivation for this modification is that the sequence of GoF $p$-values may exhibit occasional increases above the significance level $\alpha$ even when the GPD approximation is not yet appropriate. Such isolated non-rejections may lead the standard FTR rule to accept a threshold prematurely.
To reduce the impact of these random fluctuations, the robust failure-to-reject (robFTR) rule requires that the GoF test fails to reject the null hypothesis for several consecutive candidate thresholds before a threshold is accepted.

Formally, let $u_1 < u_2 < \dots < u_m$ denote the candidate thresholds and let $p_i$ denote the GoF $p$-value obtained for the GPD fit at threshold $u_i$. Let $r \ge 1$ denote a user-specified number of consecutive non-rejections required for accepting a threshold. The robust failure-to-reject rule selects the threshold
\begin{equation}
\label{eq:6perm:thresh_robftr}
u =
\min \left\{
u_i :
p_i > \alpha,\;
p_{i+1} > \alpha,\;
\dots,\;
p_{i+r-1} > \alpha
\right\},
\end{equation}
where $\alpha$ denotes the significance level of the GoF test.

Compared with the standard FTR rule (which corresponds to $r=1$), this modification reduces the probability of selecting thresholds that are accepted only due to random fluctuations in the sequence of GoF tests. Larger values of $r$ lead to more conservative threshold selection, as the GoF test must remain non-significant across a longer sequence of candidate thresholds. At the same time, excessively large values may prevent a threshold from being selected if isolated rejections occur within the candidate range.

The effect of this modification is illustrated in Figure~\ref{fig:6perm:thresh_robftr}. 
For the same Student's $t$-test permutation example considered in Section~\ref{sec:6perm:threshold:existing}, the robFTR rule selects a slightly larger threshold than the FTR method. Since the sequence of $p$-values increases smoothly after the first crossing of $\alpha$, the difference between the two thresholds remains small in this example. In practical applications, however, isolated early non-rejections may occur more frequently, in which case the robFTR rule can lead to substantially larger threshold shifts.

\begin{figure}[H]
\centering
\includegraphics[width=0.8\textwidth]{figures/permApprox/thresh_ttest_rob_ftr_ADp-1}
  \caption{\textbf{Robust failure-to-reject (robFTR) threshold selection.}
  Illustration of the same Student's $t$-test permutation example used in Section~\ref{sec:6perm:threshold:existing}. AD $p$-values obtained from fitting a GPD to exceedances above each candidate threshold. The dashed horizontal line indicates the significance level $\alpha$ of the AD test. The vertical line marks the threshold selected by the robFTR rule with $r=5$, which requires five consecutive candidate thresholds with non-rejected AD tests before accepting a threshold.
  }
\label{fig:6perm:thresh_robftr}
\end{figure}

%------------------------------------
\subsubsection*{GoF \texorpdfstring{$p$}{p}-value change point detection (GoF-CP).}
%------------------------------------

Another strategy explored in this work is based on change point detection. Instead of selecting the threshold that corresponds to the first non-rejected GoF test, this approach searches for a structural change in the sequence of AD $p$-values. The intuition behind this idea is that, for low thresholds, the GPD approximation is typically violated and the corresponding GoF tests tend to produce small $p$-values. Once the threshold enters the tail region where the GPD approximation becomes appropriate, the GoF $p$-values should increase and remain mostly above the rejection threshold. Detecting the transition between these two regimes can therefore provide an alternative criterion for threshold selection.

To identify this transition point, we apply change point detection to the sequence of AD $p$-values using the method of \citet{killick2012optimal}, as implemented in the \texttt{changepoint} \texttt{R} package \citep{killick2024changepoint}. Since the sequence of candidate thresholds may begin in a region where the GPD approximation already holds, we prepend a set of artificial small $p$-values sampled from $\text{U}(0,0.01)$ to represent an initial region of poor model fit. Change point detection is then applied to the augmented sequence of $p$-values, and the first detected change point is interpreted as the index where the goodness-of-fit behavior improves.

Let $\tilde{p}_1,\dots,\tilde{p}_{m+L}$ denote the augmented sequence of $p$-values, where the first $L$ values are artificial and the remaining values correspond to the observed AD $p$-values. If $\hat{c}$ denotes the estimated change point index, the corresponding candidate threshold is given by
\begin{equation}
u = u_{\,\min\{i : i \ge \hat{c}-L \text{ and } p_i > \alpha\}},    
\end{equation}
that is, the first candidate threshold beyond the detected change point for which the goodness-of-fit test does not reject the GPD model.

The behavior of this procedure is illustrated in Figure~\ref{fig:6perm:thresh_gofcp}. For the same permutation example considered previously, the change point method selects the same threshold as the FTR rule. The sequence of Anderson--Darling $p$-values shows a clear transition from very small values to larger values at the selected threshold. The change point detection algorithm identifies this transition and places the estimated change point at the threshold where the AD test begins to accept the GPD model. 
However, since the procedure relies on the first detected change point, it may remain sensitive to early fluctuations in the sequence of AD $p$-values. If an isolated early increase in the $p$-values occurs, the algorithm may detect this as the change point, potentially leading to premature threshold selection.

\begin{figure}[H]
\centering
\includegraphics[width=\textwidth]{figures/permApprox/thresh_ttest_gof_cp_diag-1}
  \caption{\textbf{GoF change point method.}
  Illustration of the same Student's $t$-test permutation example used in Section~\ref{sec:6perm:threshold:existing}. 
  The sequence of AD $p$-values is augmented with artificial small $p$-values (pink) to represent an initial region of poor model fit (also separated from true AD $p$-values by the vertical dotted line). The dashed horizontal line indicates the significance level $\alpha$ of the AD test. The vertical pink line marks the threshold selected by the GoF change point method.
  }
\label{fig:6perm:thresh_gofcp}
\end{figure}

%------------------------------------
\subsubsection*{Rejection proportion control (RPC)}
%------------------------------------

A further proposed threshold selection strategy exploits the same principle discussed in Section~\ref{sec:6perm:threshold:principles}: once a valid threshold has been reached, the GPD approximation should also hold for all larger thresholds. Consequently, the null hypothesis of the GoF test is expected to hold for all thresholds above the true threshold.
Under the null hypothesis, the $p$-values of the GoF tests are approximately uniformly distributed. For a significance level $\alpha$, this implies that even when the GPD model is correct, approximately a proportion $\alpha$ of the GoF tests are expected to be rejected purely by chance. Therefore, if a threshold is valid, the proportion of rejected GoF tests among all higher thresholds should not substantially exceed $\alpha$.

Based on this idea, we compute for each candidate threshold $u_i$ the proportion of rejected GoF tests among all thresholds $u_i, u_{i+1}, \dots, u_m$,
\begin{equation}
\label{eq:6perm:rpc}
r_i = \frac{1}{m-i+1} \sum_{j=i}^{m} \mathbb{I}(p_j \le \alpha),
\end{equation}
where $p_j$ denotes the GoF $p$-value at threshold $u_j$ and $\mathbb{I}(\cdot)$ is the indicator function.

To ensure that the GPD model is not rejected at the selected threshold itself, we restrict attention to thresholds with $p_i > \alpha$. 
Among these candidates, the rejection proportion control (RPC) rule selects the smallest threshold for which $r_i \le \alpha$.

If no threshold satisfies this condition, one may either not select a threshold to be conservative, or choose the threshold with the lowest rejection rate.

The behavior of this method is illustrated in Figure~\ref{fig:6perm:thresh_rpc}. 
For the same permutation example considered previously, the RPC method selects the same threshold as the FTR rule. 
The upper panel shows the sequence of AD $p$-values across candidate thresholds, while the lower panel displays the corresponding rejection proportions $r_i$. The selected threshold is the first one for which the rejection proportion among all higher thresholds falls below the significance level.

Compared with the previously discussed rules, the RPC method is less sensitive to isolated early fluctuations in the sequence of GoF $p$-values, since the decision is based on the entire sequence of tests above each candidate threshold rather than on local patterns alone. On the other hand, the method can be sensitive to many rejected GoF tests at very high thresholds, where the variability of the test statistics increases due to the small number of exceedances. Consequently, choosing the minimum number of exceedances too small may prevent the RPC method from identifying a valid threshold.

\begin{figure}[H]
\centering
\includegraphics[width=\textwidth]{figures/permApprox/thresh_ttest_rpc_ADp_diag-1.pdf}
  \caption{
  \textbf{Rejection proportion control (RPC) threshold selection.}
  The upper panel shows AD $p$-values obtained from fitting a GPD to exceedances above each candidate threshold for the same permutation example as in the previous figures. The dashed horizontal line indicates the significance level $\alpha$ of the AD test, and the vertical line marks the threshold selected by the RPC rule. The lower panel displays the proportion of rejected AD tests among all thresholds greater than or equal to each candidate threshold. The RPC method selects the smallest threshold for which this rejection proportion falls below $\alpha$, here the vertical orange line in the upper panel.
}
\label{fig:6perm:thresh_rpc}
\end{figure}

%------------------------------------
\subsubsection*{Shape-variation}
%------------------------------------

An established heuristic in extreme value analysis is to select thresholds where the estimated shape parameter $\hat{\xi}$ remains stable as the threshold increases \citep{scarrott2012review, northrop2014improved}. 
This idea is often implemented through visual inspection of so-called shape stability plots. To formalize this principle, we introduce a shape-variation rule that selects thresholds where the variability of $\hat{\xi}$ across neighboring candidate thresholds is minimal.

Let $\hat{\xi}_i$ denote the estimated shape parameter obtained from fitting the GPD at candidate threshold $u_i$. 
For each candidate threshold we compute the rolling variance of the shape estimates within a window of neighboring thresholds. 
Specifically, for a window size $w$, the variance is calculated using the estimates within a symmetric neighborhood around $u_i$, whenever sufficient observations are available.

In contrast to the GoF-based methods described above, the rolling variance is computed across all candidate thresholds, independent of the results of the GoF tests. The GoF criterion is applied only at the final selection step. 
Among all candidate thresholds for which the GoF test (we use again the AD test) is not rejected, the shape-variation rule selects the threshold with the smallest rolling variance of $\hat{\xi}$. If multiple thresholds attain the same minimum variance, the largest of these thresholds is selected.

The behavior of the method is illustrated in Figure~\ref{fig:6perm:thresh_shapevar}. 
The upper panel shows the AD $p$-values across candidate thresholds, while the lower panel displays the rolling variance of the corresponding shape parameter estimates. The selected threshold corresponds to the smallest rolling variance among thresholds for which the goodness-of-fit test is not rejected.

\begin{figure}[H]
\centering
\includegraphics[width=\textwidth]{figures/permApprox/thresh_ttest_shape_var_ADp_diag-1.pdf}
	\caption{\textbf{Shape-variation threshold selection.} The upper panel shows Anderson--Darling goodness-of-fit $p$-values obtained from fitting a GPD above each candidate threshold for the same permutation example as in the previous figures. The dashed horizontal line indicates the significance level $\alpha$ of the goodness-of-fit test. The lower panel displays the rolling variance of the corresponding estimated shape parameters. Among thresholds for which the goodness-of-fit test does not reject the GPD model (orange dots), the shape-variation rule selects the threshold with the smallest rolling variance. The selected threshold is indicated by the orange vertical line.}
\label{fig:6perm:thresh_shapevar}
\end{figure}


%====================================
\subsection{Simulation study on threshold recovery}
\label{sec:6perm:threshold:sim_thresh_recovery}
%====================================

To evaluate how well the threshold selection methods introduced in the previous section recover a true threshold, we conduct a simulation study in which the threshold is known by construction. In contrast to the previous illustrations based on permutation distributions, here the data are generated directly from parametric models that combine a bulk distribution with a GPD tail. The threshold separating these two components is specified in advance, allowing us to assess how accurately the different selection methods identify the transition between the bulk and tail regimes.

The simulation design aims to resemble the shape of typical permutation distributions encountered in practice. Such distributions often exhibit a unimodal bulk with a moderately heavy right tail corresponding to rare extreme permutation statistics. To capture this structure, we simulate data from models consisting of a bulk distribution combined with a GPD tail. Several bulk distributions are considered, including Weibull, Log-Normal, and Gamma distributions, which provide flexible models for the skewed bulk behavior commonly observed in permutation test statistics. All simulations were implemented in \texttt{R} using functions from the \texttt{evmix} package \citep{hu2018evmix}, which provides tools for simulating and fitting bulk-tail mixture models.

%------------------------------------
\subsubsection{Smooth bulk-tail transition}
%------------------------------------

As a first experiment, we consider a model in which the transition between the bulk and tail components is smooth. In this setting, the scale parameter $\sigma$ of the GPD is chosen such that the density of the combined distribution is continuous at the transition point. 
Following the construction implemented in the \texttt{evmix} package, this is achieved by choosing the GPD scale parameter according to the continuity constraint
\begin{equation}
\sigma = \frac{1 - F_W(u)}{f_W(u)},
\end{equation}
where $F_W$ and $f_W$ denote the distribution and density functions of the Weibull bulk distribution and $u$ is the nominal threshold.

Figure~\ref{fig:6perm:thresh_sim_weibull_smooth} illustrates an example of such a simulated distribution obtained from 1000 simulated observations. The upper panel shows simulated values together with the density of the generating Weibull-GPD model and the true threshold. The lower panel displays the Anderson--Darling GoF $p$-values obtained when fitting the GPD above each candidate threshold, together with the thresholds selected by the methods introduced in Sections~\ref{sec:6perm:threshold:existing} and \ref{sec:6perm:threshold:proposed}.

In this setting, all threshold selection methods identify thresholds near the beginning of the candidate threshold range and therefore far below the nominal transition point. This behavior reflects the gradual transition between the bulk and tail components: because the distribution changes smoothly at the threshold, the GPD approximation already provides a reasonable description of the upper tail for thresholds well below the nominal transition point. Consequently, the goodness-of-fit tests tend to accept the GPD model early, leading the threshold selection procedures to identify thresholds substantially smaller than the true threshold.

This example illustrates that when the transition between bulk and tail behavior is smooth, the notion of a single well-defined threshold becomes ambiguous. In such situations, goodness-of-fit based threshold selection methods may detect acceptable GPD fits well before the nominal transition between bulk and tail distributions.

\begin{figure}[H]
\centering
\includegraphics[width=\textwidth]{figures/permApprox/thresh_sim_single_weib_gpd_smooth_hist_ADp_all-1.pdf}
\caption{
  \textbf{Example of a simulated Weibull-GPD distribution with smooth transition.}
  The upper panel shows the $n=1000$ simulated observations together with the density of the generating model consisting of a Weibull bulk and a GPD tail with smoothly matched densities at the true threshold (vertical line).
  The lower panel displays Anderson--Darling (AD) $p$-values obtained when fitting a GPD to the exceedances above each candidate threshold. The dashed horizontal line indicates the significance level of the AD test. Threshold candidates are restricted so that at least 10 exceedances remain for fitting, which causes the candidate threshold range to end before the right boundary of the histogram. Colored markers indicate the thresholds selected by the different threshold detection methods.
}
\label{fig:6perm:thresh_sim_weibull_smooth}
\end{figure}

%------------------------------------
\subsubsection{Sharp bulk-tail transition}
%------------------------------------

Since the smooth construction considered above is not well-suited for assessing the methods’ ability to locate the true threshold, we next consider models in which the GPD scale parameter is not chosen to enforce continuity at the threshold. Instead, the parameters of the bulk distribution and the GPD tail are specified directly so that the resulting distribution exhibits a visible, but not overly extreme, jump at the transition point. This construction yields a threshold that is sufficiently pronounced to be detectable by the threshold selection procedures, while still producing distributions that resemble realistic test statistic distributions with a dominant bulk and a moderately heavy right tail.

To examine the robustness of the threshold selection methods across different bulk shapes, we consider three bulk distributions: Weibull, Log-Normal, and Gamma distributions. Each of these is combined with a GPD tail with shape parameter $\xi=-0.2$, representing a moderate choice within the bounded-tail regime ($\xi<0$), which is the primary focus of this work. The parameters of the bulk and tail components are chosen such that the transition at the threshold is neither completely smooth nor unrealistically abrupt, thereby providing the methods with a genuine opportunity to recover the true threshold. The complete parameter settings used in the simulations are summarized in Table~\ref{tab:threshold_sim_settings}. In each scenario, a sample of size $n=1000$ is simulated from the corresponding bulk-tail model.

\begin{table}[t]
\centering
\setlength{\tabcolsep}{12pt}
\caption{Parameter settings used in the threshold recovery simulations.}
\label{tab:threshold_sim_settings}
\begin{tabular}{lccc}
\hline
Scenario & Bulk parameters & GPD parameters & Start exceedances \\
\hline
Weibull-GPD   & $k=1.2,\ \lambda=1$      & $\xi=-0.2,\ \sigma=1.8$ & 250, 1000 \\
Lognormal-GPD & $\mu=0,\ \sigma=0.8$     & $\xi=-0.2,\ \sigma=2.5$ & 250, 1000 \\
Gamma-GPD     & $k=1.2,\ \theta=0.6$     & $\xi=-0.2,\ \sigma=1.3$ & 250, 1000 \\
\hline
\end{tabular}
\end{table}

Figure~\ref{fig:6perm:thresh_sim_four_methods} illustrates example realizations from these models. Panels (a)-(c) show simulated data from the Weibull-GPD, Log-normal-GPD, and Gamma-GPD models, respectively. 
In each case, the upper plot displays the simulated observations together with the density of the generating model and the true threshold, while the lower plot shows the AD $p$-values obtained when fitting a GPD to the exceedances above each candidate threshold. The colored markers indicate the thresholds selected by the different detection methods.

In panels (a)-(c), the threshold search starts with the smallest observed samples ($k_0=1000$), to evaluate the behavior of the methods across the complete range of observations. In practice, however, threshold selection is typically restricted to the upper part of the distribution in order to avoid fitting extreme value models to the bulk region. 
Panel (d) illustrates this situation using the Gamma-GPD model. Here, the threshold search is restricted to the upper 25\% of the sample distribution by excluding the first 750 observations from the candidate threshold range ($k_0=250$). This setup more closely reflects practical applications, where threshold detection is typically performed only within the upper fraction of the distribution.

\begin{figure}[H]
\centering
\includegraphics[width=\textwidth]{figures/permApprox/thresh_sim_single_combined_four_panels-1.pdf}
\caption{
  \textbf{Examples of simulated distributions with a sharp transition between bulk and tail components.}
  \textbf{(a)-(c):} Weibull-GPD, Log-normal-GPD, and Gamma-GPD models with $k_0=1000$ starting exceedances.
  \textbf{(d):} Gamma-GPD model with restricted threshold search ($k_0=250$).
  In each panel, the upper plot shows the $n=1000$ simulated observations together with the density of the generating model and the true threshold (vertical line). The lower plot displays Anderson--Darling (AD) $p$-values obtained when fitting a GPD to the exceedances above each candidate threshold. The dashed horizontal line indicates the significance level of the AD test. Threshold candidates are restricted so that at least 10 exceedances remain for fitting, which causes the candidate threshold range to end before the right boundary of the histogram. Colored markers indicate the thresholds selected by the different threshold detection methods.
}
\label{fig:6perm:thresh_sim_four_methods}
\end{figure}

%------------------------------------
\subsubsection{Simulation results with repeated samples}
%------------------------------------

The examples in Figure~\ref{fig:6perm:thresh_sim_four_methods} illustrate the behavior of the different threshold selection methods for individual simulated datasets. As already indicated, the shape of the AD $p$-value sequence can vary substantially between realizations, in particular in the presence of local fluctuations below the true threshold. Consequently, the selected threshold may depend on the specific random sample and thus on the simulation seed.

To assess the performance of the methods more systematically, we conducted a simulation study with repeated sampling. For each setting in Table~\ref{tab:threshold_sim_settings}, we generated 100 independent samples of size $n=1000$ and applied all threshold selection methods. Performance was evaluated in terms of the deviation of the estimated threshold $\hat{u}$ from the true threshold $u_{\text{true}}$.

Figure~\ref{fig:6perm:thresh_sim_rep_1000exc} shows the results for $k_0 = 1000$, i.e., when the threshold search starts at the smallest observation. It is immediately apparent that the RPC method is the only approach that is not sensitive to early increases in the AD $p$-value sequence. This robustness arises because the method aggregates information across all candidate thresholds rather than relying on local behavior. 
However, this comes at the cost of increased variability, with RPC exhibiting a comparatively wide spread of estimated thresholds across all settings. Similarly, the ForwardStop and Shape-Var methods show substantial variability and occasional large deviations from the true threshold.

In contrast, the FTR and robFTR methods, as well as the GoF-CP approach, exhibit low variability and estimates that are consistently close to the true threshold across all considered distributions. While these methods are more sensitive to early increases in the AD $p$-values, this limitation is less critical in practice, as threshold search is typically not performed starting from the smallest observation.

\begin{figure}[H]
\centering
\includegraphics[width=\textwidth]{figures/permApprox/thresh_sim_rep_jitter_box_all_exc1000_step1-1.pdf}
\caption{
  \textbf{Threshold estimation accuracy for $k_0 = 1000$.}
  Distribution of estimation errors $\hat{u} - u_{\text{true}}$ across 100 simulated samples for Weibull-GPD, Log-normal-GPD, and Gamma-GPD models. 
  The threshold search starts at the smallest observation, i.e., all observations are considered as potential exceedances. 
  Points represent individual simulation runs and boxplots summarize the distribution across repetitions. 
}
\label{fig:6perm:thresh_sim_rep_1000exc}
\end{figure}

Figure~\ref{fig:6perm:thresh_sim_rep_250exc} presents the corresponding results for $k_0 = 250$, where the threshold search is restricted to the upper part of the distribution. Overall, the results are similar to the $k_0 = 1000$ case, but without the pronounced outliers caused by early acceptances of the GoF tests. In addition, the variability of most methods is reduced. The FTR and robFTR methods, together with GoF-CP, continue to show the most stable performance with estimates close to the true threshold, while the remaining methods still exhibit larger variability.

\begin{figure}[H]
\centering
\includegraphics[width=\textwidth]{figures/permApprox/thresh_sim_rep_jitter_box_all_exc250_step1-1.pdf}
\caption{
  \textbf{Threshold estimation accuracy for $k_0 = 250$.}
  Distribution of estimation errors $\hat{u} - u_{\text{true}}$ across 100 simulated samples for Weibull-GPD, Log-normal-GPD, and Gamma-GPD models. 
  The threshold search is restricted to the upper part of the distribution by considering only the largest 250 observations as candidate exceedances. 
  Points represent individual simulation runs and boxplots summarize the distribution across repetitions. 
}
\label{fig:6perm:thresh_sim_rep_250exc}
\end{figure}

%====================================
\subsection{Simulation study on \texorpdfstring{$p$}{p}-value recovery}
\label{sec:6perm:threshold:sim_pval_recovery}
%====================================

While the previous subsection focused on the ability of the different methods to recover the true threshold in controlled settings, the primary objective of this work is the accurate approximation of permutation-based $p$-values. It is therefore natural to evaluate the threshold selection methods also in terms of their impact on $p$-value approximation, which is done in this section.

To study the effect of threshold selection on $p$-value approximation under controlled conditions, we consider a two-sample \textit{t}-test with Gaussian data. 
This setting serves as a benchmark because the parametric \textit{t}-test is correctly specified under the data-generating model and therefore provides a natural reference for assessing the quality of the approximated permutation $p$-values.

We generated two groups of normally distributed observations with equal group sizes $n \in \{25,50,75,100,250,500,750,1000\}$ and standard deviation $\sigma=1$. For each sample size, we considered effect sizes $d \in \{0.5,1,1.5,2\}$, corresponding to the absolute mean difference between the two groups. 

The test statistic is given by the two-sample \textit{t}-statistic under the equal-variance assumption,
\begin{equation}
t = \frac{\bar{x}_1 - \bar{x}_2}{\sqrt{(s_1^2 + s_2^2)/n}}
\end{equation}
where $\bar{x}_1, \bar{x}_2$ denote the sample means and $s_1^2, s_2^2$ the sample variances of the two groups. 
For equal group sizes, this expression is equivalent to the classical pooled-variance \textit{t}-statistic.
For every configuration, $B=1000$ random label permutations were used to generate the permutation distribution of the test statistic. A total of 1000 independent repetitions were performed per setting in order to assess variability.

Permutation $p$-values were approximated using the constrained GPD framework described in Section~\ref{sec:6perm:gpd_tail_modeling}. To isolate the effect of threshold selection, all other components of the procedure were fixed throughout this comparison: GPD parameters were estimated using the LME estimator, the support constraint was imposed via the SLLS $\varepsilon$-rule, and the threshold search was initialized at $k_0=250$ exceedances.

All threshold selection methods introduced in Sections~\ref{sec:6perm:threshold:existing} and \ref{sec:6perm:threshold:proposed} are evaluated. For RPC, we distinguish between two strategies when no threshold satisfies the rejection proportion criterion: the standard RPC approach, which does not return a threshold in this case, and a modified variant (RPC-Min) that instead selects the threshold with the smallest rejection proportion. For the FTR-based methods (FTR and robFTR), candidate thresholds are evaluated on a grid with step size 10, whereas for the remaining methods (FwdStop, GoF-CP, RPC, RPC-Min, Shape-Var), all admissible thresholds above $k_0$ are considered, as coarser grids may adversely affect their performance.

Figure~\ref{fig:6perm:thresh_pvalue_recovery} summarizes the results. 
The left panel shows the ratio between the approximated permutation $p$-values and the reference $p$-values obtained from the parametric $t$-test, while the right panel displays the number of exceedances selected by each method.

The results reveal that the two classical approaches (FTR and FwdStop) as well as the proposed methods robFTR and GoF-CP yield highly comparable $p$-value approximations. Across all sample sizes, the ratios to the reference \textit{t}-test $p$-values are close to one and exhibit similar variability. This similarity is also reflected in the selected thresholds: for these methods, the number of exceedances is close to the initial value $k_0$ in most cases, and the variability in threshold selection remains low.

The Shape-Var method exhibits a different behavior. While its median $p$-value ratio is slightly closer to 1, it more frequently produces ratios below one, indicating an underestimation of the \textit{t}-test $p$-values, which is undesirable in this context. This behavior is explained by the substantially larger thresholds selected by Shape-Var, resulting in smaller numbers of exceedances and, consequently, increased variability. This is clearly visible in the right panel, where both the level and variability of the selected thresholds differ markedly from the other methods.

The two variants of the RPC method show a distinct pattern. 
For the standard RPC approach, the rejection proportion often does not fall below the $\alpha$ level, so that no valid threshold is identified. In these cases, the method falls back to the empirical $p$-value, leading to extremely large ratios. 
The modified RPC-Min variant avoids this issue by always selecting a threshold, but at the cost of increased variability in the selected thresholds. Moreover, it yields $p$-values below the reference more often than the other methods.

\begin{figure}[H]
\centering
\includegraphics[width=\textwidth]{figures/permApprox/sim_ttest_thresh_combined_pvals_nexceed_by_n-1.pdf}
\caption{
  \textbf{Comparison of threshold selection methods for permutation $p$-value approximation.}
  \textbf{Left:} Ratios of approximated permutation $p$-values to the reference $p$-values obtained from the parametric \textit{t}-test ($p_{\text{method}} / p_{\text{\textit{t}-test}}$). 
  \textbf{Right:} Number of exceedances selected by each method. 
  Results are shown for sample sizes $n \in \{50, 100, 500, 1000\}$ based on 1000 simulation repetitions per setting. 
  Permutation distributions are generated using $B=1000$ permutations, and threshold selection is initialized at $k_0 = 250$ exceedances. 
  The methods compared include FTR, robFTR, FwdStop, GoF-CP, Shape-Var, and two variants of rejection proportion control: RPC and RPC-Min.
}
\label{fig:6perm:thresh_pvalue_recovery}
\end{figure}

%====================================
\subsection{Study on the number of starting exceedances}
\label{sec:6perm:threshold:sim_exceed0}
%====================================

In addition to the choice of the threshold detection method, the number of starting exceedances $k_0$ plays a crucial role in the GPD-based approximation framework. 
As discussed in Section~\ref{sec:6perm:threshold:principles}, the threshold search is typically restricted to the upper part of the permutation distribution in order to avoid fitting the GPD to the bulk. 
The parameter $k_0$ determines how many of the largest permutation statistics are considered as candidates for threshold selection and thus directly influences both the bias and variability of the resulting $p$-value estimates.

To investigate the impact of $k_0$, we consider the same simulation setting as described in Section~\ref{sec:6perm:threshold:sim_pval_recovery}. 
In brief, two samples are drawn from Gaussian distributions and compared using a two-sample \textit{t}-test, with permutation distributions generated from $B \in \{500, 1000, 5000, 10000\}$ permutations. The constrained GPD approximation is applied using the LME estimator with support constraint based on the SLLS $\epsilon$ rule. Here, we use the FTR method for threshold detection.
We compare five choices for $k_0$: two fixed values ($k_0 = 100$ and $k_0 = 250$) as well as three relative fractions of the permutation size ($k_0 = 0.1B$, $0.25B$, and $0.5B$). 

Figure~\ref{fig:6perm:exceed0} summarizes the results. 
The figure shows the ratio between the approximated permutation $p$-values and the reference $p$-values obtained from the parametric \textit{t}-test across different values of $B$ and $k_0$.
The results reveal a clear dependence of the $p$-value approximation accuracy on the choice of $k_0$, which varies with the number of permutations $B$. 
For $B = 500$, using $k_0 = 250$ (corresponding to 50\% of $B$) results in more conservative $p$-value approximations with substantially lower variation than smaller values of $k_0$.  
For $B = 1000$, both $k_0 = 250$ (=$0.25B$) and $k_0 = 500$ yield acceptable variability, with median ratios closest to one for $k_0 = 250$.  
For larger $B$, the variability decreases up to $k_0 = 0.25B$ and increases thereafter, while the median remains closest to the \textit{t}-test reference at $k_0 = 0.25B$.  

\begin{figure}[H]
\centering
\includegraphics[width=\textwidth]{figures/permApprox/sim_ttest_nexc_pval_ratios_by_B-1.pdf}
\caption{
  \textbf{Influence of the number of starting exceedances $k_0$ on $p$-value approximation.}
  Ratios of approximated permutation $p$-values to the reference \textit{t}-test $p$-values ($p_{\text{method}} / p_{\text{\textit{t}-test}}$) for different numbers of permutations $B$. 
  Each panel corresponds to a fixed value of $B$, and boxplots summarize results across 1000 simulation repetitions. 
  The $x$-axis shows different choices of $k_0$, including fixed values and fractions of $B$. The constrained GPD approximation is performed using LME estimation and FTR-based threshold detection.
  The stars mark our recommended choice ($k_0 = \mathrm{max}(0.25B, 250)$).
  This figure is adapted from \citep{peschel2026permapprox}.
}
\label{fig:6perm:exceed0}
\end{figure}

%====================================
\subsection{Default threshold detection approach}
\label{sec:6perm:threshold:decision}
%====================================

The simulation studies presented in this section jointly address two key components of threshold selection in the \texttt{permApprox} workflow: the choice of the threshold detection method and the number of starting exceedances $k_0$. 

Across both threshold recovery and $p$-value recovery settings, the robFTR method consistently provided the most favorable balance between accuracy, stability, and computational efficiency (Sections~\ref{sec:6perm:threshold:sim_thresh_recovery} and \ref{sec:6perm:threshold:sim_pval_recovery}). In particular, it yielded $p$-value approximations that were close to the reference values while exhibiting reduced variability compared to alternative methods. Moreover, as a variant of the FTR approach, it allows for early stopping once a suitable threshold is identified, thereby improving computational efficiency. Compared to the standard FTR approach, the robust version reduces sensitivity to local fluctuations in the goodness-of-fit statistics and thus mitigates the risk of premature threshold selection.

The study on the number of starting exceedances (Section~\ref{sec:6perm:threshold:sim_exceed0}) showed that the choice of $k_0$ substantially affects the bias-variance trade-off of the resulting $p$-value estimates. Across all considered permutation sizes, a relative choice of $k_0 = 0.25B$ provided the most stable and accurate results, while fixed values such as $k_0 = 250$ were preferable for smaller numbers of permutations.

Based on these findings, we adopt the combination of robFTR for threshold detection and $k_0 = \max(0.25B, 250)$ as the default configuration in the \texttt{permApprox} workflow.

%%%%%%%%%%%%%%%%%%%%%%%%%%%%%%%%%%%%%
\section{Derivation of a data-adaptive \texorpdfstring{$\varepsilon$}{epsilon} rule}
\label{sec:6perm:epsilon}
%%%%%%%%%%%%%%%%%%%%%%%%%%%%%%%%%%%%%

%====================================
\subsection{Motivation}
\label{sec:6perm:epsilon:motivation}
%====================================

While the constraint introduced in Section~\ref{sec:6perm:constraint} itself guarantees that the resulting tail probability is well-defined, the choice of the safety margin $\varepsilon$ has a substantial influence on the resulting $p$-value approximation.
Conceptually, $\varepsilon$ controls how far the fitted support boundary must extend beyond the observed excess. 
If $\varepsilon$ is chosen too small, the evaluation point may lie extremely close to the fitted upper endpoint of the GPD. In this situation, the resulting tail probability $\bar F_{\mathrm{GPD}}$ becomes numerically unstable and may still collapse to zero due to machine precision limitations. Conversely, choosing $\varepsilon$ too large pushes the support boundary far into the upper tail, which leads to overly conservative $p$-values. 

Importantly, the optimal magnitude of $\varepsilon$ cannot be determined independently of the scale and variability of the test statistic. The distance between the observed statistic and the bulk of the permutation distribution depends on factors such as the sample size, the variability of the permutation statistics, and the underlying test statistic itself. A fixed constant $\varepsilon$ on the original scale of the test statistic would therefore lead to inconsistent behavior across different tests and experimental settings.

These considerations motivate the development of a data-adaptive rule for selecting $\varepsilon$ that satisfies several design requirements:

\begin{itemize}
  \item \textbf{Scale invariance.}  
  The rule should behave consistently across test statistics with different magnitudes and units.
  \item \textbf{Numerical stability.}  
  The safety margin must ensure that the evaluation point remains sufficiently separated from the fitted support boundary to avoid numerical underflow or instability in tail probability evaluation.
  \item \textbf{Smooth adaptation to tail extremeness.}  
  The rule should allow the safety margin to increase as the observed statistic moves further into the tail of the permutation distribution, while avoiding excessive conservativeness for moderate statistics.
  \item \textbf{Robustness across test statistics.}  
  The rule should remain applicable across different test statistics and distributional settings rather than relying on properties of a specific test.
\end{itemize}

%====================================
\subsection{Simulation setup}
\label{sec:6perm:epsilon:sim_setup}
%====================================

The design requirements formulated in the previous subsection imply that the safety margin $\varepsilon$ cannot be chosen as a fixed constant but must depend on the observed test statistic and the properties of the permutation distribution. 
However, the precise functional form of such a rule is not immediately obvious. To derive a suitable construction, we therefore conducted a series of simulation experiments to systematically investigate how different choices of $\varepsilon$ influence the resulting $p$-value approximation.

As a benchmark scenario, we considered the two-sample Student's $t$-test. 
Data were generated from two independent Gaussian populations with equal variance,
\begin{equation}
    X_i \sim \mathcal{N}(\mu,1), \qquad Y_i \sim \mathcal{N}(0,1),
\end{equation}
where the parameter $\mu$ determines the effect size. 
To obtain test statistics covering a broad range of tail positions, we considered effect sizes $d \in \{0,\,0.5,\,1,\,1.5,\,2\}$.
For each effect size, $200$ independent tests were generated, resulting in a total of $1000$ tests per simulation setting. 
For every dataset, the one-sided two-sample $t$-test statistic was computed to test
\begin{equation}
    H_0:\mu_X - \mu_Y = 0 
    \qquad \text{versus} \qquad 
    H_1:\mu_X - \mu_Y > 0 .
\end{equation}
To investigate the influence of the sample size on the behavior of the $\varepsilon$ rules, we repeated this experiment for group sample sizes $n \in \{50,\,100,\,250,\,500,\,750,\,1000\}$.
For each value of $n$, the simulation setting therefore consists of $1000$ tests covering the five effect sizes considered above.

Permutation testing was used to approximate the null distribution of the statistic. 
For each test, a permutation distribution was generated using $B=1000$ randomly sampled permutations of the group labels. 
These permutation distributions served as the basis for the GPD-based tail approximation under different choices of the safety margin $\varepsilon$.

To evaluate the behavior of candidate $\varepsilon$ rules, the resulting permutation-based p-value approximations were compared with reference p-values obtained from the classical Student's $t$-test. Since the Gaussian data-generating model satisfies the assumptions of the $t$-test, these reference values provide a reliable benchmark for assessing the accuracy of the tail approximation.

This simulation setup produces observed test statistics ranging from values close to the bulk of the permutation distribution to extreme values located far in the tail. 
It therefore provides a suitable environment for studying how the safety margin $\varepsilon$ should depend on the extremeness of the observed statistic.

%====================================
\subsection{From linear dependence to the log-saturation rule}
\label{sec:6perm:epsilon:log_saturation}
%====================================

Using the simulation setup described above, we now derive a suitable $\varepsilon$ rule for the two-sample $t$-test. In the following, we describe the successive refinements that led from simple heuristic constructions to the final standardized lifted log-saturation (SLLS) rule presented in Section~\ref{sec:6perm:epsilon:final_rule}.

%------------------------------------
\subsubsection*{Dependence on the test statistic} 
%------------------------------------

A natural first approach is to relate the safety margin directly to the magnitude of the observed test statistic. 
For each test $j$, we consider the safety margin $\varepsilon_j$ as a function of the corresponding observed statistic $T_{\mathrm{obs}}^{(j)}$. 
The simplest construction of this type assumes a linear dependence,
\begin{equation}
\varepsilon_j = c\,\big|T_{\mathrm{obs}}^{(j)}\big|,
\end{equation}
where $c>0$ is a constant scaling factor.

Figure~\ref{fig:6perm:eps_rule_linear} illustrates this construction for several sample sizes. 
The approximated $p$-values obtained from the constrained GPD approximation are shown together with the corresponding $t$-test $p$-values. 
The linear rule stabilizes the approximation across the range of observed statistics, confirming that the safety margin must indeed increase with $|T_{\mathrm{obs}}|$. 
However, the figure also shows that no single constant $c$ yields good agreement with the $t$-test $p$-values across all considered sample sizes. 

\begin{figure}[H]
  \centering
  \includegraphics[width=\linewidth]{figures/permApprox/eps_rule_factor01_3n.pdf}
  \caption{\textbf{Linear $\varepsilon$ rule.}
  Constrained GPD $p$-values obtained using the rule $\varepsilon = 0.1\,|T_{\mathrm{obs}}|$ for Gaussian two-sample $t$-tests with different sample sizes ($n=50,100,500$). Colored points show the approximated $p$-values, while open circles indicate the corresponding $p$-values from the classical $t$-test. 
  Colors represent the underlying effect sizes $d \in \{0,0.5,1,1.5,2\}$. 
  While the linear rule stabilizes the approximation as the observed statistic increases, the deviations across sample sizes indicate that the rule must incorporate an explicit dependence on $n$.
}
  \label{fig:6perm:eps_rule_linear}
\end{figure}

%------------------------------------
\subsubsection*{Incorporating sample size}
%------------------------------------

The results obtained with the linear rule indicate that the safety margin must depend not only on the magnitude of the observed statistic but also on the sample size. 
In particular, Figure~\ref{fig:6perm:eps_rule_linear} shows that the resulting approximations become increasingly conservative as the sample size grows. 
This suggests that the safety margin should decrease with increasing sample size, indicating an inverse relationship between $\varepsilon$ and $n$.

To account for this effect, we considered constructions of the form
\begin{equation}
  \varepsilon_j = \big|T_{\mathrm{obs}}^{(j)}\big| \frac{c}{n^p},
\end{equation}
where $n$ denotes the per-group sample size and $c>0$ and $p>0$ are tuning parameters. 
This formulation reduces the magnitude of the safety margin as the sample size increases and was therefore expected to counteract the overly conservative behavior observed for the purely linear rule.

Figure~\ref{fig:6perm:eps_rule_n_squared} illustrates the resulting $p$-value approximations for three representative sample sizes ($n=50,100,500$). Across a range of parameter choices, the setting $c=1000$ and $p=2$ yielded $p$-value approximations closest to the corresponding $t$-test $p$-values and is shown in the figure. Compared with the purely linear rule, incorporating the factor $n^{-2}$ largely removes the dependence of the approximation on the sample size, leading to much more consistent behavior across different values of $n$.

However, the dependence of $\varepsilon$ on $|T_{\mathrm{obs}}|$ remains purely linear. 
Consequently, the resulting approximations still fail to capture the pronounced curvature in the relationship between the $t$-statistic and its corresponding $p$-value. This observation motivates the introduction of nonlinear rules for the safety margin in the next step of the derivation.

\begin{figure}[H]
  \centering
  \includegraphics[width=\linewidth]{figures/permApprox/eps_rule_n_squared_3n.pdf}
  \caption{\textbf{Sample-size scaled linear $\varepsilon$ rule.}
  Constrained GPD $p$-values obtained using the rule $\varepsilon = |T_{\mathrm{obs}}|\,1000/n^2$ for Gaussian two-sample $t$-tests with sample sizes $n=50$, $100$, and $500$. 
  Colored points show the approximated $p$-values, while open circles indicate the corresponding $p$-values from the classical $t$-test. 
  Colors represent the true effect sizes $d \in \{0,0.5,1,1.5,2\}$. 
  Incorporating the factor $n^{-2}$ substantially reduces the dependence of the approximation on the sample size. 
  Nevertheless, the linear dependence on $|T_{\mathrm{obs}}|$ remains unable to reproduce the nonlinear relationship between the $t$-statistic and its $p$-value.
  }
  \label{fig:6perm:eps_rule_n_squared}
\end{figure}

%------------------------------------
\subsubsection*{Logarithmic growth and saturation}
%------------------------------------

A natural choice to accommodate the observed curvature of the $t$-test $p$-values is to replace the linear dependence on the observed statistic by a logarithmic transformation. 
In our exploratory experiments, we considered rules of the form
\begin{equation}
\varepsilon_j = \log\!\left(1 + c\,\big|T_{\mathrm{obs}}^{(j)}\big|\right)\frac{1}{n^p},
\end{equation}
where the factor $n^{-p}$ preserves the sample-size scaling introduced in the previous subsection. 
Compared with the linear rules, this formulation increases more rapidly for small to moderate statistics while flattening for larger values, thereby reflecting the nonlinear relationship between the test statistic and its corresponding $p$-value.

However, the logarithmic function still grows without bound as $\lvert T_{\mathrm{obs}}^{(j)} \rvert$ increases. 
To avoid an unnecessarily large safety margin for extreme statistics, it is therefore desirable to normalize the logarithmic growth so that the resulting function saturates at a finite upper level.

To this end, let $T_{\max}=\max_j \big|T_{\mathrm{obs}}^{(j)}\big|$ denote the largest observed statistic across all tests, representing the upper end of the relevant range of test statistics. 
We then consider a saturating logarithmic rule of the form
\begin{equation}
\varepsilon_j =
\varepsilon_{\max}(n)\,
\frac{\log\!\left(1 + \kappa(n)\,\big|T_{\mathrm{obs}}^{(j)}\big|\right)}
     {\log\!\left(1 + \kappa(n)\,T_{\max}\right)} ,
\label{eq:eps_log_sat}
\end{equation}
where $\varepsilon_{\max}(n)$ determines the maximal safety margin reached when $\lvert T_{\mathrm{obs}}^{(j)} \rvert = T_{\max}$, and $\kappa(n)>0$ controls the curvature of the logarithmic growth.

The numerator describes a logarithmic increase of the safety margin with increasing extremeness of the observed statistic, while the normalization in the denominator ensures that the rule saturates smoothly at the plateau level $\varepsilon_{\max}(n)$ as $\lvert T_{\mathrm{obs}}^{(j)} \rvert$ approaches $T_{\max}$. 
The parameter $\kappa(n)$ determines how rapidly the curve bends toward saturation: larger values produce a steeper initial increase and earlier saturation, whereas smaller values yield a more gradual growth.

Instead of introducing the sample-size dependence through an explicit factor of the form $n^{-p}$, the dependence on the sample size is now incorporated into the parameters $\varepsilon_{\max}(n)$ and $\kappa(n)$ of the saturation rule. 
We therefore parameterize these quantities as
\begin{equation}
\kappa(n) = \kappa_{\mathrm{factor}}\frac{n_{\mathrm{ref}}}{n},
\qquad
\varepsilon_{\max}(n) = \tau\,\frac{n_{\mathrm{ref}}}{n},
\label{eq:eps_scaling}
\end{equation}
where $n$ denotes the per-group sample size and $n_{\mathrm{ref}}$ is a fixed reference sample size. This formulation preserves the empirical observation that the safety margin should decrease as the sample size increases, while allowing the curvature and the plateau height of the logarithmic rule to adapt simultaneously to $n$.

The reference value $n_{\mathrm{ref}}$ defines a pivot point at which the strength of the regularization transitions from stronger regularization for smaller samples to weaker regularization for larger samples. 
Our simulation studies indicated that $n_{\mathrm{ref}}=500$ provides $p$-value approximations close to the corresponding $t$-test $p$-values across a wide range of scenarios and is therefore recommended as a general default.

The constants $\kappa_{\mathrm{factor}}$ and $\tau$ control the curvature and the plateau height of the logarithmic saturation rule, respectively. 
Specifically, $\kappa_{\mathrm{factor}}$ determines how rapidly the rule bends toward its saturation level, whereas $\tau$ determines the magnitude of the maximal safety margin $\varepsilon_{\max}(n)$. 
Our simulation studies showed that the parameter choices $\kappa_{\mathrm{factor}}=1000$ and $\tau=0.25$ consistently yielded $p$-value approximations that were close to the corresponding $t$-test reference values across a broad range of sample sizes and effect sizes. We therefore recommend these values as default settings for the rule.

To assess the behavior of the resulting rule, Figure~\ref{fig:6perm:eps_rule_log_sat} shows constrained GPD $p$-values obtained with the logarithmic saturation rule for three representative sample sizes.

The upper panels (a) display the approximated $p$-values as a function of the observed statistic. Compared with the linear constructions depicted in Figure~\ref{fig:6perm:eps_rule_n_squared}, the logarithmic saturation rule captures the curvature of the $t$-test $p$-values much more accurately across the full range of test statistics. The lower panels (b) compare the approximated $p$-values directly with the corresponding $t$-test $p$-values.

Overall, the approximated $p$-values follow the $t$-test reference values closely, confirming that the logarithmic saturation rule substantially improves the approximation of the permutation tail. However, for test statistics in an intermediate range the approximated $p$-values tend to fall slightly below the corresponding $t$-test values. 
In the lower panels, this behavior appears as points located above the dashed identity line. The effect becomes particularly pronounced for larger sample sizes, as visible in the $n=500$ panel.

Since our goal is to avoid underestimation of the reference $p$-values, this observation indicates that an additional correction is required for moderate test statistics.

\begin{figure}[H]
  \centering
  \includegraphics[width=\linewidth]{figures/permApprox/eps_rule_log_sat_pp.pdf}
  \caption{
  \textbf{Logarithmic saturation $\varepsilon$ rule.}
  Results for Gaussian two-sample $t$-tests with sample sizes $n=50$, $100$, and $500$. 
  \textbf{(a)} $p$-values plotted against the observed statistic $|t_{\mathrm{obs}}|$. 
  Colored dots represent constrained GPD approximations and open black circles represent the corresponding $t$-test $p$-values. 
  \textbf{(b)} Comparison of approximated and $t$-test $p$-values, where the colored dots represent p-values and the dashed line denotes the identity. 
  In all panels, colors indicate the effect size ($d \in \{0,0.5,1,1.5,2\}$).
  Dashed circles indicate the region where a correction is needed.
  The $p$-values are shown on a $\log_{10}$ scale with tick labels in the original $p$-value scale.
  }
  \label{fig:6perm:eps_rule_log_sat}
\end{figure}

%====================================
\subsection{Lift for intermediate statistics}
\label{sec:6perm:epsilon:lift}
%====================================

To correct the localized deviation for intermediate $p$-values observed in Figure~\ref{fig:6perm:eps_rule_log_sat}, while preserving the desirable behavior for extreme statistics, we introduce an additional lift component to the safety margin.

The lift is defined using the compactly supported Wendland $C^2$ kernel from the $d=3$ Wendland family \citep{wendland1995piecewise},
\begin{equation}
\psi(l) = (1-l)^4(1+4l), \qquad l\in[0,1],
\label{eq:6perm:wendland}
\end{equation}
which is smooth, nonnegative, and vanishes together with its first two derivatives at $l=1$. 
Compactly supported Wendland functions are widely used as smoothing kernels in mesh-free particle methods such as smoothed particle hydrodynamics (SPH) and in simulations of fluid-structure interaction, where their finite support enables stable local interpolation \citep{beckert2001multivariate,liu2010smoothed}.
These properties ensure that the lift increases the safety margin primarily for small to moderate values of $l_j$, while fading out smoothly near the cap and leaving the saturation behavior for extreme statistics unchanged.

To apply the lift on a unified scale, we express the position of the observed statistic relative to the range of observed statistics as
\begin{equation}
\label{eq:6perm:rel_position}
l_j = \frac{\big|T_{\mathrm{obs}}^{(j)}\big|}{T_{\max}} \in [0,1].
\end{equation}
The Wendland lift term is then applied as a smooth function of this relative position.

Combining the logarithmic saturation rule from Equation~\eqref{eq:eps_log_sat} with the lift component yields the modified safety margin
\begin{equation}
\varepsilon_j =
\varepsilon_{\max}(n)\,
\frac{\log\!\left(1 + \kappa(n)\,\big|T_{\mathrm{obs}}^{(j)}\big|\right)}
     {\log\!\left(1 + \kappa(n)\,T_{\max}\right)}
\;+\;
\rho_{\mathrm{lift}}\,\psi(l_j),
\label{eq:6perm:epsilon_lls}
\end{equation}
where $\rho_{\mathrm{lift}}>0$ controls the magnitude of the localized lift. In the experiments considered in this work, the value $\rho_{\mathrm{lift}}=3$ yielded accurate approximations and is therefore recommended as a default.

Figure~\ref{fig:6perm:eps_rule_lls} summarizes the effect of the lift. 
The upper panels (a) shows the resulting $\varepsilon$ functions as a function of the observed statistic. 
Compared with the purely logarithmic rule, the lifted version introduces a localized increase of the safety margin for intermediate statistics while leaving the behavior for very small and very large values of $|T_{\mathrm{obs}}|$ largely unchanged. 

The middle and lower panels (b) and (c) illustrate the resulting constrained GPD $p$-values and their comparison with the corresponding $t$-test $p$-values. 
With the lift component, the approximated $p$-values now closely follow the curvature of the $t$-test reference values across the entire range of statistics. 
For larger sample sizes, the approximated $p$-values tend to lie slightly above the $t$-test values, resulting in a mildly conservative behavior, which is preferable to systematically producing smaller values.

\begin{figure}[H]
  \centering
  \includegraphics[width=\linewidth]{figures/permApprox/eps_rule_lls.pdf}
  \caption{
  \textbf{Lifted log-saturation $\varepsilon$ rule.}
  Results for Gaussian two-sample $t$-tests with sample sizes $n=50$, $100$, and $500$. 
  \textbf{(a)} $\varepsilon$-rules as functions of the observed statistic $|t_{\mathrm{obs}}|$, comparing the lifted log-saturation rule with the logarithmic saturation rule. 
  \textbf{(b)} Resulting $p$-values plotted against the observed statistic $|t_{\mathrm{obs}}|$, where colored dots represent constrained GPD approximations and open black circles denote the corresponding $t$-test $p$-values. 
  \textbf{(c)} Comparison of approximated and $t$-test $p$-values, where colors indicate the effect size ($d \in \{0,0.5,1,1.5,2\}$) and the dashed line denotes the identity. 
  The $p$-values are shown on a $\log_{10}$ scale with tick labels in the original $p$-value scale.
  }
  \label{fig:6perm:eps_rule_lls}
\end{figure}

%====================================
\subsection{Permutation-based cap}
\label{sec:6perm:epsilon:cap}
%====================================

The lifted log-saturation rule presented in Equation~\eqref{eq:6perm:epsilon_lls} normalizes the observed statistic using the largest observed value,
$T_{\max} = \max_j \big|T_{\mathrm{obs}}^{(j)}\big|$,
which defines the upper end of the statistic range used in the normalization. 
While this choice is natural in large-scale multiple testing settings, where many observed statistics are available, it can become sensitive to extreme outlying statistics and may be unstable when only a small number of tests is considered.

To improve robustness, we replace $T_{\max}$ by a more general \emph{cap} value $T_{\mathrm{cap}}$ representing the upper range of plausible test statistics. 
This cap can be defined either from the observed statistics or from the permutation distribution.

\noindent\textbf{Observed cap.}
When many tests are available, the cap can be defined using the largest observed statistic,
\begin{equation}
    T_{\mathrm{cap}} = T_{\max}.
\end{equation}
This choice exploits the empirical range of the observed statistics and is therefore appropriate in large-scale multiple testing scenarios.

\noindent\textbf{Permutation-based cap.}
In settings with few tests or when robustness against extreme observed statistics is desired, the cap can instead be derived from the permutation distribution. 
Specifically, we define
\begin{equation}
T_{\mathrm{cap}} = Q_q\!\big(|\mathcal{T}|\big)\, g(n),
\end{equation}
where $Q_q(|\mathcal{T}|)$ denotes a high quantile (e.g., $q=0.99$) of the absolute permutation statistics and $g(n) = \max\{1,\alpha\sqrt{n}\}$ is a mild inflation factor depending on the sample size $n$. This inflation prevents the cap from becoming too small when the permutation distribution exhibits limited variability.

Using either construction, the relative position of the observed statistic \eqref{eq:6perm:rel_position} is changes to
\begin{equation}
\label{eq:6perm:rel_position_cap}
l_j = \frac{\big|T_{\mathrm{obs}}^{(j)}\big|}{T_{\mathrm{cap}}} \in [0,1].
\end{equation}

Figure~\ref{fig:6perm:eps_rule_lls_perm} shows that in the Gaussian simulation setting the permutation-based cap leads to $p$-value approximations that are almost indistinguishable from those obtained with the maximum-based cap. This behavior can be explained by the design of the simulation, which includes a wide range of effect sizes and therefore produces a broad spectrum of observed test statistics. In such an idealized setting, the largest observed statistic already provides a reliable estimate of the upper range of plausible statistics, so replacing $T_{\max}$ by the permutation-based cap has little practical effect. The resulting $p$-values are slightly smaller than those obtained with the maximum-based cap, reflecting the somewhat smaller $\varepsilon$ values produced by the permutation-based normalization. Nevertheless, the approximations remain generally slightly larger than the corresponding $t$-test $p$-values, preserving the conservative behavior observed for the lifted log-saturation rule.

\begin{figure}[H]
  \centering
  \includegraphics[width=\linewidth]{figures/permApprox/eps_rule_lls_perm.pdf}
  \caption{\textbf{Lifted log-saturation rule with permutation-based cap.}
  Constrained GPD $p$-values obtained using the lifted log-saturation rule with permutation-based cap for Gaussian two-sample $t$-tests with sample sizes $n=50$, $100$, and $500$. 
  Colored points represent the constrained GPD approximations, while open black circles denote the corresponding $t$-test $p$-values. 
  Colors indicate the effect size ($d \in \{0,0.5,1,1.5,2\}$). 
  The $p$-values are shown on a $\log_{10}$ scale with tick labels in the original $p$-value scale.
  }
  \label{fig:6perm:eps_rule_lls_perm}
\end{figure}

%====================================
\subsection{Standardization}
\label{sec:6perm:epsilon:standardization}
%====================================

The lifted log-saturation rule derived in the previous subsections was motivated and calibrated using the two-sample $t$-test. 
In this setting, the test statistic is approximately standardized and symmetric, providing a natural and stable scale on which the $\varepsilon$-rule can operate. 
However, many commonly used test statistics differ substantially in their distributional properties and numerical scale. 
To ensure that the $\varepsilon$-rule remains applicable across a broad range of statistical tests, it is therefore necessary to examine its behavior in settings where the test statistic deviates strongly from the approximately symmetric and standardized structure of the $t$-statistic.

To this end, we repeated the simulation analysis using the nonparametric Wilcoxon rank-sum test based on the Mann-Whitney $U$ statistic. 
This test provides a particularly challenging scenario for the $\varepsilon$-rule. 
First, the distribution of the $U$ statistic is typically highly skewed under alternative hypotheses, making it difficult to approximate by simple parametric models. 
Second, the numerical range of the statistic differs drastically from that of the $t$-statistic. 
While $t$-values typically lie within a relatively small range, the $U$ statistic grows on the order of $O(n^2)$ and can therefore take values in the thousands even for moderate sample sizes. 

To investigate whether the lifted log-saturation rule remains stable under such conditions, we simulated two groups of exponentially distributed observations with means $1$ and $1+d$, corresponding to effect sizes
$
d \in \{0,\,0.5,\,1,\,1.5,\,2\}.
$
For each effect size, $200$ independent tests were generated, resulting in $1000$ tests per simulation setting. 
Permutation distributions with $B=1000$ replicates were used for the constrained GPD-based approximations, while the classical Wilcoxon test provided reference $p$-values.

Since the Wilcoxon test is nonparametric, there is no simple parametric expression for the exact tail probability of the statistic. In the following experiments we therefore use the classical Wilcoxon p-values obtained from the asymptotic normal approximation of the Mann-Whitney statistic (without continuity correction) as reference values.

Figure~\ref{fig:6perm:eps_comparison_wilcox} illustrates the effect of applying the lifted log-saturation rule to the Wilcoxon setting. 
The left panel shows the constrained GPD $p$-values obtained when the rule is applied directly to the raw Mann-Whitney $U$ statistics. In this case, the resulting approximations deviate substantially from the Wilcoxon reference $p$-values for nonzero effect sizes. 

\begin{figure}[H]
  \centering
  \includegraphics[width=\linewidth]{figures/permApprox/eps_rule_wilcox.pdf}
  \caption{\textbf{Effect of standardization on constrained permutation $p$-values for the Wilcoxon test.}
  Data were generated from exponential distributions with means $1$ and $1+d$ ($d \in \{0,0.5,1,1.5,2\}$), using equal group sizes ($n=100$) and $200$ tests per effect size ($1000$ in total).
  Colored points show constrained GPD $p$-values, while open circles indicate reference $p$-values obtained from the asymptotic Wilcoxon rank-sum test (normal approximation without continuity correction).
  The left panel applies the lifted log-saturation rule directly to the raw Mann-Whitney $U$ statistics, whereas the right panel uses the standardized lifted log-saturation rule based on standardized $Z$ scores.
  This figure is adapted from \citep{peschel2026permapprox}.
  }
  \label{fig:6perm:eps_comparison_wilcox}
\end{figure}

To address this issue, we derive the $\varepsilon$-rule on a standardized scale and subsequently map it back to the scale of the original test statistic. 
For test $j$, let
\begin{equation}
Z_{\mathrm{obs}}^{(j)} =
\frac{\lvert T_{\mathrm{obs}}^{(j)} - \mu_j \rvert}{\sigma_j}
\end{equation}
denote the standardized observed statistic, where $\mu_j$ and $\sigma_j$ are the mean and standard deviation of the corresponding permutation distribution. 
The lifted log-saturation rule is then formulated on this standardized scale rather than directly on the raw test statistic.

To localize the lift on a common scale, we again express the relative position of the observed statistic within the relevant range of values. 
On the standardized scale, Equation \eqref{eq:6perm:rel_position_cap} becomes
\begin{equation}
\label{eq:6perm:rel_position_std}
l_j = \frac{Z_{\mathrm{obs}}^{(j)}}{Z_{\mathrm{cap}}} \in [0,1],
\end{equation}
where $Z_{\mathrm{cap}}>0$ denotes a standardized cap defining the upper end of the relevant statistic range. 
Analogously to the construction in Section~\ref{sec:6perm:epsilon:cap}, this cap can either be based on the largest observed standardized statistic across tests or derived from a high quantile of the standardized permutation distribution.

The lifted log-saturation rule is then applied as a function of the standardized statistic $Z_{\mathrm{obs}}^{(j)}$ and the relative position $l_j$, yielding a standardized safety margin $\varepsilon_j^\ast$. 
Finally, the safety margin is mapped back to the scale of the original statistic by
\begin{equation}
\varepsilon_j = \sigma_j \, \varepsilon_j^\ast ,
\end{equation}
so that the constrained GPD approximation remains defined on the original test statistic scale.

The right panel of Figure~\ref{fig:6perm:eps_comparison_wilcox} shows the resulting constrained GPD $p$-values obtained using this standardized construction. 
In contrast to the rule applied directly to the raw $U$ statistic, the standardized version produces approximations that closely follow the Wilcoxon reference $p$-values across all considered effect sizes. 
Although the Wilcoxon $p$-values themselves are obtained from the asymptotic normal approximation of the rank-sum statistic rather than from an exact parametric model, they provide a useful reference for assessing the stability of the approximation. 
These results therefore demonstrate that deriving the $\varepsilon$-rule on a standardized scale is essential for obtaining robust behavior across test statistics with different numerical ranges and distributional properties.

%====================================
\subsection{Final \texorpdfstring{$\varepsilon$}{epsilon} rule}
\label{sec:6perm:epsilon:final_rule}
%====================================

The previous subsections derived the $\varepsilon$ rule in a stepwise manner. 
Starting from simple linear constructions, we introduced a logarithmic saturation component, added a localized lift to correct deviations for intermediate statistics, and derived the rule on a standardized scale to ensure robustness across test statistics with different numerical ranges. 
Combining these elements yields the \emph{standardized lifted log-saturation} (SLLS) rule.

Let $l_j$ denote the relative position of the standardized observed statistic as specified in Equation~\eqref{eq:6perm:rel_position_std}. The standardized safety margin is then defined as
\begin{equation}
\varepsilon_j^{\ast}
=
\max\!\left\{
\varepsilon^{\ast}_{\max}(n)\,
\frac{\log\!\big(1+\kappa(n)\,l_j\big)}{\log\!\big(1+\kappa(n)\big)}
\;+\;
\rho_{\mathrm{lift}}\,
\psi(l_j)
\;,\;
\varepsilon_{\min}^{\ast}
\right\}.
\label{eq:epsilon_standardized}
\end{equation}

The first term in Equation~\eqref{eq:epsilon_standardized} corresponds to the logarithmic saturation rule derived in Section~\ref{sec:6perm:epsilon:log_saturation}, which ensures a smooth increase of the safety margin with increasing extremeness of the observed statistic while saturating at a finite upper level. 
The second term represents the localized lift that corrects the systematic deviations for intermediate statistics introduced in Section~\ref{sec:6perm:epsilon:lift}, where $\psi(l)$ denotes the Wendland lift function. 
The outer maximum ensures a strictly positive lower bound $\varepsilon_{\min}^{\ast}>0$, preventing the safety margin from becoming too small and thereby avoiding numerical instability in the tail evaluation.

Both the curvature and the plateau height of the saturation component depend on the sample size. 
Specifically, we set
\begin{equation}
\kappa(n)=\kappa_{\mathrm{factor}}\frac{n_{\mathrm{ref}}}{n},
\qquad
\varepsilon^\ast_{\max}(n)=\tau\,\frac{n_{\mathrm{ref}}}{n},
\label{eq:epsilon_scaling}
\end{equation}
where $n$ denotes the per-group sample size and $n_{\mathrm{ref}}$ is a fixed reference sample size (Section~\ref{sec:6perm:epsilon}). 
When group sizes differ, we use $n=\min(n_1,n_2)$ to obtain a conservative scaling.

Finally, the safety margin on the original scale of the test statistic is obtained by rescaling with the permutation standard deviation via $\varepsilon_j=\sigma_j\,\varepsilon_j^\ast$.

Unless stated otherwise, the default parameter values used in this work are
\[
n_{\mathrm{ref}}=500,
\qquad
\kappa_{\mathrm{factor}}=1000,
\qquad
\tau=0.25,
\qquad
\rho_{\mathrm{lift}}=3,
\qquad
\varepsilon_{\min}^{\ast}=10^{-6}.
\]

%=========================================
\subsection{Data-adaptive \texorpdfstring{$\varepsilon$}{epsilon} refinement}
\label{sec:6perm:epsilon:refinement}
%=========================================

Even with the SLLS $\varepsilon$-rule, machine-underflow $p$-values (reported as zero) may still occur for negative shape estimates ($\hat\xi < 0$), either because the test statistic is extremely large or because the fitted upper support is very close to the constraint, causing the tail probability to fall below machine precision. We therefore employ a \textit{data-adaptive refinement} for the target factor $\tau$, one of the SLLS tuning parameters, to identify the minimum $\tau$ value that eliminates underflow zeros while preserving moderate cases.

Let $\varepsilon(t_{\text{obs}};\tau)$ denote the SLLS epsilon produced at target factor $\tau$ and define the index set of underflowing tests at $\tau$ as

\begin{equation}
 \mathcal Z(\tau)=\{j:\ \hat p_j(\tau)\leq \text{.Machine.double.xmin}\}.
\end{equation}

We seek the minimal $\tau^\star=\inf \{\tau:\ |\mathcal Z(\tau)|=0\}$.
Assuming monotonicity of the “no-underflow” event with increasing $\tau$ (empirically observed), we use a two-phase search:

\begin{itemize}
    \item \textbf{Phase 1 (expansion / bracketing):} Starting at $\tau_0$ (default 0.25), we step additively by $s$ and inflate the step geometrically $s\leftarrow s\cdot\rho$ until $|\mathcal Z(\tau)|=0$ is achieved, producing a bracket $[\tau_{\text{lo}},\tau_{\text{hi}}]$ with underflow at $\tau_{\text{lo}}$ and none at $\tau_{\text{hi}}$. To reduce cost, at each try we fit the GPD only for tests currently in $\mathcal Z(\tau)$.

    \item \textbf{Phase 2 (bisection / refinement):} Bisection of $[\tau_{\text{lo}},\tau_{\text{hi}}]$ until the width falls below a numerical tolerance, returning $\tau^\star\approx\tau_{\text{hi}}$. During bisection we again refit only the currently underflowing subset. Finally, we perform one full refit for all selected tests at $\tau^\star$ and report the resulting $p$-values.
\end{itemize}

\subsubsection*{Defaults and stopping}
Expansion step $s_{\text{init}}=10$, growth $\rho=\frac{1+\sqrt5}{2}\approx1.618$, at most 20 expansion steps; bisection tolerance 0.1 with at most 30 iterations. The procedure logs $\tau$ and $|\mathcal Z(\tau)|$ per iteration for transparency. If expansion hits its cap with $|\mathcal Z(\tau)|>0$, we return the best $\tau$ found and flag this in the messages (users can increase \texttt{step\_init}, \texttt{grow}, or the iteration caps).

\subsubsection*{Computational notes}
The search is subset-aware: only tests with underflow are refit during exploration, yielding substantial speedups when zeros are rare. The final results are obtained from a single full refit at $\tau^\star$, so intermediate exploration does not affect the reported estimates. Progress bars (optional) mirror those used in threshold detection.

This refinement aims to eliminate, up to machine precision, all reported zeros due to underflow while keeping $\tau$ as small as possible, thereby preserving power and avoiding overly conservative inflation of $\varepsilon$ in routine cases. If the expansion phase hits its iteration cap with remaining underflowing tests, the best available $\tau$ is returned and this is flagged in the output.



%%%%%%%%%%%%%%%%%%%%%%%%%%%%%%%%%%%%%
\section{GPD parameter estimation}
\label{sec:6perm:param_est}
%%%%%%%%%%%%%%%%%%%%%%%%%%%%%%%%%%%%%

The previous sections established the key components required for GPD-based tail modeling of permutation distributions. In particular, we introduced a support-constrained formulation of the GPD, derived a data-adaptive rule for the evaluation point, and discussed strategies for selecting an appropriate threshold to define the tail region. Once a threshold has been selected, the remaining task is to estimate the parameters of the Generalized Pareto distribution from the resulting exceedances.

Accurate parameter estimation is crucial for reliable $p$-value approximation, as both bias and variability in the estimated shape and scale parameters directly affect the extrapolation into the extreme tail. At the same time, this setting poses several challenges. 
First, the number of exceedances is often limited, especially when the threshold is chosen conservatively. Second, the estimation is performed under a support constraint to ensure valid $p$-values, which may violate standard assumptions of classical estimation procedures. 
These aspects motivate a careful comparison of different estimation approaches in the context of permutation testing.

In the following, we introduce a range of established and recently proposed methods for estimating the GPD parameters. These methods are compared in the subsequent simulation studies, where we evaluate their ability to recover the true GPD parameters under both correctly specified models and under misspecification induced by the support constraint.

%====================================
\subsection{Methods for GPD parameter estimation}
\label{sec:6perm:param_est:methods}
%====================================

Let $Y_1, \dots, Y_k$ denote the excesses over a given threshold $u$, defined as $Y_i = X_i - u$ for all observations $X_i > u$. 
These excesses are assumed to follow a GPD as defined in Section~\ref{sec:6perm:gpd_tail_modeling} with shape parameter $\xi \in \mathbb{R}$ and scale parameter $\sigma > 0$. We consider several estimation methods that have been proposed in the extreme value literature. Each method provides estimates $(\hat{\xi}, \hat{\sigma})$ based on the observed excesses, but differs in how these estimates are obtained.

In addition to the unconstrained parameter estimates $(\hat{\xi}, \hat{\sigma})$, we consider constrained estimates $(\hat{\xi}_{(C)}, \hat{\sigma}_{(C)})$ obtained by enforcing the support condition introduced in Section~\ref{sec:6perm:gpd_tail_modeling}. 
Specifically, for $\xi < 0$, the parameter space is restricted to values satisfying
\begin{equation}
-\frac{\sigma}{\xi} > Y_{\mathrm{obs}} + \varepsilon,
\end{equation}
ensuring that the fitted distribution is defined at the evaluation point. 
As a consequence, the constrained estimates induce a boundary $\hat{s}_{(C)}=-\hat{\sigma}_{(C)}/\hat{\xi}_{(C)}$ satisfying $\hat{s}_{(C)} > Y_{\mathrm{obs}} + \varepsilon$.
This constraint restricts the admissible parameter space and may introduce bias if the true distribution violates the imposed support condition, but is necessary to ensure valid evaluation of tail probabilities.

In the following, we briefly introduce the considered estimation methods and describe how each method is adapted to obtain constrained parameter estimates.

%------------------------------------
\subsubsection*{Method of moments (MOM)}
%------------------------------------

Introduced by \citet{hosking1987parameter}, the method of moments (MOM) estimates the GPD parameters by matching the theoretical moments to the empirical moments of the excesses. 
Let $\bar{Y}$ and $s^2$ denote the sample mean and variance of the excesses. 
The MOM estimators are given by
\begin{equation}
\hat{\xi} = \frac{1}{2} \left( \frac{\bar{Y}^2}{s^2} - 1 \right), 
\qquad
\hat{\sigma} = \frac{1}{2} \bar{Y} \left( \frac{\bar{Y}^2}{s^2} + 1 \right),
\end{equation}
given the moments exist.

The method of moments is simple and computationally efficient but does not involve optimization and cannot incorporate support constraints. We include it only for comparison in the unconstrained setting.

%------------------------------------
\subsubsection*{Two-dimensional MLE (MLE2D)}
%------------------------------------

The classical two-parameter maximum likelihood estimation (MLE2D) \citep{pickands1975statistical} maximizes the log-likelihood over both shape and scale. It is widely used in extreme value theory and serves as a standard reference estimator.

Given excesses $Y_1, \dots, Y_k$, the log-likelihood function is
\begin{equation}
\ell(\xi, \sigma) = -k \log \sigma - \left(1 + \frac{1}{\xi}\right) \sum_{i=1}^k \log\left(1 + \frac{\xi Y_i}{\sigma}\right),
\end{equation}
where the parameter space is restricted by $1+\xi Y_i/\sigma>0$ for all $i$, ensuring that all observed excesses lie in the support of the fitted GPD. The estimates $(\hat{\xi}, \hat{\sigma})$ are obtained by numerically maximizing this likelihood.

Since the optimization is sensitive to the choice of starting values, we initialize the parameter search by probability-weighted moments (PWM) \citep{greenwood1979probability}. In the constrained setting, the boundary point $Y_{\mathrm{obs}}+\varepsilon$ is included when computing these starting values in order to obtain an initialization that is compatible with the restricted parameter space. If the resulting PWM estimates are degenerate or infeasible, the method of moments is used as a fallback. This strategy was found to yield substantially more stable optimization behavior than MOM-based initialization alone, particularly in constrained settings where the unrestricted MOM estimate may suggest an excessively negative shape parameter.

In the constrained setting, the optimization is performed over the restricted parameter space defined by
$-\frac{\sigma}{\xi} > Y_{\mathrm{obs}}+\varepsilon$ for $\xi < 0$.
In practice, this condition is incorporated directly into the likelihood evaluation by treating parameter configurations that violate the constraint as infeasible and assigning a penalty to the objective function. As a result, the constrained estimates $(\hat{\xi}_{(C)}, \hat{\sigma}_{(C)})$ induce a boundary satisfying $\hat{s}_{(C)} > Y_{\mathrm{obs}}+\varepsilon$, ensuring that the fitted distribution is defined at the evaluation point.

%------------------------------------
\subsubsection*{One-dimensional MLE (MLE1D)}
%------------------------------------

Castillo and Serra \citep{castillo2015likelihood} propose a reparametrization of the GPD in terms of the endpoint of the distribution. 
In their notation, the parameters $(\kappa, \psi, \sigma)$ correspond to our $(\xi, \sigma, s)$, where $s = -\sigma/\xi$ denotes the upper endpoint in the bounded-tail case ($\xi < 0$). 

Using this parametrization, the likelihood can be expressed as a function of the endpoint $s$ alone. 
In our notation, the corresponding shape parameter is given by
\begin{equation}
\xi(s) = \frac{1}{k} \sum_{i=1}^k \log\left(1 - \frac{Y_i}{s}\right),
\end{equation}
and the scale parameter follows as $\sigma(s) = -\xi(s)\,s$. 
Substituting these expressions into the log-likelihood yields a profile likelihood depending only on $s$, which can be optimized numerically. 
This reduces the estimation problem to a one-dimensional optimization and directly targets the parameter that determines the support of the distribution.

As in the two-dimensional MLE approach, the constrained estimation is implemented by penalizing parameter configurations that violate the constraint $s_{(C)} > Y_{\mathrm{obs}} + \varepsilon$ in the case $\xi < 0$. 

%------------------------------------
\subsubsection*{Likelihood moment estimation (LME)}
%------------------------------------

The likelihood moment estimation (LME) method proposed by \citet{zhang2007likelihood} combines likelihood-based and moment-based ideas within a one-dimensional estimation framework. In the original formulation, the GPD is parametrized by scale and shape parameters $\sigma$ and $k$, where $k = -\xi$ in our notation.

The method is based on the reparametrization $b = -\xi/\sigma$, so that $\sigma = -\xi/b$. For the excesses $Y_1,\dots,Y_k$, the log-likelihood can be expressed in terms of $b$, after expressing $\xi$ as a function of $b$ via the likelihood equations, as
\begin{equation}
\ell(b) = k \log\left(\frac{b}{-\xi}\right) + \left(1 - \frac{1}{\xi}\right)\sum_{i=1}^k \log(1 - bY_i),
\end{equation}
with $b < 1/Y_{(k)}$ in the bounded-support case.

From the likelihood equations, the shape parameter can be written as a function of $b$,
\begin{equation}
\xi(b) = \frac{1}{k}\sum_{i=1}^k \log(1 - bY_i),
\end{equation}
and consequently $\sigma(b) = -\xi(b)/b$. Thus, the original two-parameter estimation problem reduces to estimating the single parameter $b$.

Instead of solving the corresponding one-dimensional likelihood equation directly, the LME approach replaces it by the estimating equation
\begin{equation}
\frac{1}{k}\sum_{i=1}^k (1 - bY_i)^p = \frac{1}{1 - r},
\end{equation}
where $r < 1$ is a tuning parameter and
\begin{equation}
p = \frac{r}{\frac{1}{k}\sum_{i=1}^k \log(1 - bY_i)}.
\end{equation}
This construction links moment and likelihood estimation: specific choices of $r$ recover moment-type equations, while the likelihood equations arise as limiting cases. In our implementation, we use the default choice $r = -1/2$ as suggested by \citet{zhang2007likelihood}.

Once $\hat{b}$ is obtained, the parameter estimates follow as
\begin{equation}
\hat{\xi} = \frac{1}{k}\sum_{i=1}^k \log(1 - \hat{b}Y_i),
\qquad
\hat{\sigma} = -\frac{\hat{\xi}}{\hat{b}}.
\end{equation}

In the constrained setting, the support requirement is incorporated directly into the iterative solution for $b$. More precisely, the upper bound governing the admissible range of $b$ is not taken as the largest observed excess alone, but as
$Y_{\max} = \max\{Y_{(k)},\, Y_{\mathrm{obs}}+\varepsilon\}$ and the admissible range of $b$ is restricted to values satisfying $b < Y_{\max}^{-1}$.

%------------------------------------
\subsubsection*{Zhang-Stephens estimator (ZSE)}
%------------------------------------

The estimator proposed by \citet{zhang2009new} extends the likelihood moment estimation approach by incorporating a grid-based stabilization inspired by Bayesian ideas. While the LME method determines the parameter $b = -\xi/\sigma$ as the solution of a nonlinear estimating equation, the Zhang-Stephens estimator instead constructs $b$ as a weighted average over a set of candidate values.

As in the LME approach, the likelihood can be expressed in terms of the single parameter $b$ by profiling out $\xi$, yielding a profile likelihood $L(b)$. Rather than maximizing this likelihood or solving an estimating equation, the ZSE defines the estimator
\begin{equation}
\hat{b} = \frac{\int b \, \pi(b) L(b)\, db}{\int \pi(b) L(b)\, db},
\end{equation}
where $\pi(b)$ is a data-dependent weighting function.

In practice, this expression is approximated by a finite weighted sum over a grid of candidate values $b_1,\dots,b_m$ within the admissible parameter range. The corresponding weights are given by normalized likelihood values,
\begin{equation}
w_j = \frac{L(b_j)}{\sum_{\ell=1}^m L(b_\ell)},
\end{equation}
where the effect of the weighting function $\pi(b)$ is implicitly reflected in the construction of the grid. The estimator is then given by
\begin{equation}
\hat{b} = \sum_{j=1}^m b_j w_j.
\end{equation}

This construction replaces the root-finding problem of LME by a weighted averaging procedure and thereby reduces sensitivity to irregularities of the likelihood surface. The method can be interpreted as a discretized empirical Bayes estimator, where the weighting function $\pi(b)$ induces a stabilization of the likelihood, particularly in moderate sample sizes.

Once $\hat{b}$ is obtained, the parameter estimates follow as
\begin{equation}
\hat{\xi} = \frac{1}{k}\sum_{i=1}^k \log(1 - \hat{b}Y_i),
\qquad
\hat{\sigma} = -\frac{\hat{\xi}}{\hat{b}}.
\end{equation}

In the constrained setting, the support requirement is incorporated through the construction of the candidate grid. Specifically, the admissible range of $b$ is determined by
$Y_{\max} = \max\{Y_{(k)},\, Y_{\mathrm{obs}}+\varepsilon\},$
and the grid is restricted to values satisfying $b_j < Y_{\max}^{-1}$. Since $\hat{b}$ is a weighted average of admissible values, the resulting constrained estimates $(\hat{\xi}_{(C)}, \hat{\sigma}_{(C)})$ satisfy $\hat{s}_{(C)} > Y_{\mathrm{obs}}+\varepsilon$ whenever $\xi < 0$.

%------------------------------------
\subsubsection*{Nonlinear least squares estimator (NLS2)}
%------------------------------------

The nonlinear least squares estimator proposed by \citet{song2012quantile} estimates the GPD parameters by minimizing the deviations between the empirical distribution function (EDF) of the excesses $Y_1,\dots,Y_k$ and the corresponding theoretical GPD distribution function. Let $F_{\mathrm{emp}}$ denote the EDF of the excesses and $F_{\xi,\sigma}$ the corresponding GPD distribution function. The estimator is defined as
\begin{equation}
(\hat{\xi}, \hat{\sigma}) = \arg\min_{\xi,\sigma} \sum_{i=1}^k \left(F_{\mathrm{emp}}(Y_i) - F_{\xi,\sigma}(Y_i)\right)^2.
\end{equation}

A direct minimization of this objective function is numerically unstable, as the GPD distribution function is highly sensitive to the shape parameter $\xi$. In particular, small deviations in $\xi$ can lead to large changes in $F_{\xi,\sigma}$, resulting in a poorly behaved optimization problem.

To address this issue, \citet{song2012quantile} propose a two-step fitting procedure (NLS2). In the first step, the problem is transformed using the identity
\begin{equation}
\log\left(1 - F_{\xi,\sigma}(x)\right) = -\frac{1}{\xi} \log\left(1 + \frac{\xi x}{\sigma}\right).
\end{equation}
Defining
\begin{equation}
y_i = \log\left(1 - F_{\mathrm{emp}}(Y_i)\right),
\end{equation}
the transformed model can be written as
\begin{equation}
y_i = -\frac{1}{\xi} \log\left(1 + \frac{\xi Y_i}{\sigma}\right) + \varepsilon_i,
\end{equation}
and the parameters are estimated by minimizing the corresponding residual sum of squares via nonlinear least squares.

In the second step, the resulting estimates are used as initial values for minimizing the original least squares criterion. This two-stage procedure improves numerical stability and yields more reliable estimates compared to direct nonlinear least squares fitting.

In contrast to likelihood-based approaches such as MLE, LME, and ZSE, the NLS2 estimator is based on matching empirical and theoretical distribution functions rather than maximizing a likelihood. This reduces sensitivity to numerical instability of the likelihood surface, but typically results in lower efficiency when the GPD model is correctly specified.

In the constrained setting, the support requirement is incorporated by penalizing parameter values for which the fitted GPD is not defined at the relevant boundary point. More precisely, the effective boundary is taken as
$Y_{\max} = \max\{Y_{(k)},\, Y_{\mathrm{obs}}+\varepsilon\}$,
and both steps of the NLS2 procedure are carried out only over parameter values for which the fitted distribution remains valid at $Y_{\max}$. For $\xi < 0$, this corresponds to requiring
$-\frac{\sigma}{\xi} > Y_{\max}$.
In the implementation, parameter values violating this condition are assigned a large penalty in the objective function. As a result, the constrained estimates $(\hat{\xi}_{(C)}, \hat{\sigma}_{(C)})$ satisfy $-\hat{\sigma}_{(C)}/\hat{\xi}_{(C)} > Y_{\max}$.

%------------------------------------
\subsubsection*{Weighted nonlinear least squares likelihood moments estimator (WNLLSM)}
%------------------------------------

The weighted nonlinear least squares likelihood moments estimator (WNLLSM) proposed by \citet{zhao2019new} combines ideas from likelihood-based profiling and weighted least squares fitting of the distribution function. It is introduced as a refinement of the weighted nonlinear least squares moments estimator (WNLSM), with the goal of improving estimation accuracy by eliminating the shape parameter from the optimization problem. In our setting, the method is applied directly to the excesses $Y_1,\dots,Y_k$.

As in the likelihood-based estimators above, the method is based on the reparametrization
$b = \xi/\sigma$, so that the shape parameter can be written as a function of $b$,
\begin{equation}
\xi(b) = \frac{1}{k}\sum_{i=1}^k \log(1 + bY_i),
\end{equation}
with the support condition $1 + bY_i > 0$ for all $i$. This reduces the estimation problem to a one-dimensional optimization in $b$.

The estimator is constructed in two steps. In the first step, a preliminary estimate of $b$ is obtained by minimizing an unweighted least squares criterion based on the transformed tail distribution. In the second step, this criterion is refined by introducing weights
\begin{equation}
\omega_i = \frac{(k+1)^2(k+2)}{i(k-i+1)},
\end{equation}
which are motivated by the variance of the order statistics of a uniform distribution. Larger order statistics therefore receive larger weights, so that the estimation focuses more strongly on the upper tail. The final estimator is obtained as
\begin{equation}
\hat{b}
=
\arg\min_b
\sum_{i=1}^k
\omega_i
\left(
\frac{i}{k+1}
-
F_{\xi(b),b}(Y_{(i)})
\right)^2,
\end{equation}
where $F_{\xi(b),b}$ denotes the GPD distribution function written in terms of the parameters $(\xi,b)$ and $Y_{(i)}$ are the $i$-th order statistic of the excesses.

Once $\hat{b}$ is obtained, the remaining parameters follow as
\begin{equation}
\hat{\xi} = \frac{1}{k}\sum_{i=1}^k \log(1 + \hat{b}Y_i),
\qquad
\hat{\sigma} = \frac{\hat{\xi}}{\hat{b}}.
\end{equation}

Compared to WNLSM, the WNLLSM estimator avoids direct joint optimization over both shape and scale parameters by profiling out $\xi$ and optimizing only over $b$. This yields a simpler and typically more stable estimation procedure, while preserving the emphasis on upper-order observations through the weights. In this sense, the method combines a likelihood-based reduction of the parameter space with a weighted distribution-matching criterion. 

In the constrained setting, the support requirement is incorporated directly into the one-dimensional optimization for $b$. More precisely, the effective boundary is taken as
\begin{equation}
Y_{\max} = \max\{Y_{(k)},\, Y_{\mathrm{obs}}+\varepsilon\},
\end{equation}
and the admissible range of $b$ is restricted to values satisfying $b > -1/Y_{\max}$.
For $\xi < 0$, this is equivalent to requiring that the fitted support boundary satisfies
$-\frac{\sigma}{\xi} > Y_{\max}$.
In the implementation, this constraint is enforced through the lower bound of the one-dimensional optimization.

%====================================
\subsection{Simulation study on parameter recovery}
\label{sec:6perm:param_est:sim_unconstrained}
%====================================

We begin by evaluating the ability of the considered estimators to recover the true parameters of the Generalized Pareto distribution in an idealized setting without support constraints. 
This setting serves as a baseline for comparing the intrinsic statistical efficiency of the estimators, independent of the additional hurdles introduced by the constraint required for valid $p$-value approximation.

To this end, exceedances were generated directly from a GPD with known parameters,
\begin{equation}
Y_i \sim \mathrm{GPD}(\sigma_0 = 1,\, \xi_0),
\qquad
\xi_0 \in \{-0.6, -0.4, -0.2, 0, 0.2, 0.4, 0.6\},
\end{equation}
covering both bounded ($\xi_0 < 0$) and heavy-tailed ($\xi_0 > 0$) regimes. 
Sample sizes were varied as
%\begin{equation}
$n \in \{25, 50, 75, 100, 250, 500, 750, 1000\}$,
%\end{equation}
and for each parameter configuration, $R = 500$ independent samples were generated. 
For each replicate, the parameters $(\hat{\sigma}, \hat{\xi})$ were estimated using the methods introduced in Section~\ref{sec:6perm:param_est:methods} via the \texttt{fit\_gpd()} function of the \texttt{permApprox} package.

Estimation accuracy was quantified using the empirical root mean squared error (RMSE) across simulation replicates:,
\begin{equation}
\mathrm{RMSE}(\xi) = 
\sqrt{\frac{1}{R}\sum_{r=1}^{R}(\hat{\xi}_r - \xi_0)^2},
\qquad
\mathrm{RMSE}(\sigma) = 
\sqrt{\frac{1}{R}\sum_{r=1}^{R}(\hat{\sigma}_r - \sigma_0)^2},
\end{equation}
where lower RMSE values indicate higher estimation accuracy.

Figures~\ref{fig:6perm:sim_unconstrained_rmse_shape} and~\ref{fig:6perm:sim_unconstrained_rmse_scale} summarize the RMSE for the shape and scale parameters, respectively, across all estimators, sample sizes, and true shape values. 
The color scale represents absolute RMSE values, enabling direct comparison of estimation accuracy across methods and parameter settings. 
In addition, Figure~\ref{fig:6perm:sim_unconstrained_error_shape} shows the distribution of absolute estimation errors for $\xi_0 = -0.4$ across sample sizes.

\begin{figure}[H]
  \centering
  \includegraphics[width=\textwidth]{figures/permApprox/sim_GPD_unconstr_heat_shape_rmse-1.pdf}
  \caption{Root mean squared error (RMSE) of the estimated GPD shape parameter $\hat{\xi}$ across sample sizes $n$, estimation methods, and true shape values $\xi_0\in \{-0.6, -0.4, -0.2, 0.2, 0.4, 0.6\}$. 
  The color scale represents absolute RMSE values, ranging from the minimum to the maximum observed across all scenarios. Lighter colors correspond to lower RMSE and thus higher estimation accuracy. The starred method (LME*) corresponds to the default estimator in \texttt{permApprox}.
  This figure is adapted from \citep{peschel2026permapprox}.}
  \label{fig:6perm:sim_unconstrained_rmse_shape}
\end{figure}

\begin{figure}[H]
  \centering
  \includegraphics[width=\textwidth]{figures/permApprox/sim_GPD_unconstr_heat_scale_rmse-1.pdf}
  \caption{Root mean squared error (RMSE) of the estimated GPD scale parameter $\hat{\sigma}$ across sample sizes $n$, estimation methods, and true shape values $\xi_0\in \{-0.6, -0.4, -0.2, 0.2, 0.4, 0.6\}$. 
  The color scale represents absolute RMSE values, ranging from the minimum to the maximum observed across all scenarios. Lighter colors correspond to lower RMSE and thus higher estimation accuracy. The starred method (LME*) corresponds to the default estimator in \texttt{permApprox}. 
  This figure is adapted from \citep{peschel2026permapprox}.}
  \label{fig:6perm:sim_unconstrained_rmse_scale}
\end{figure}

\begin{figure}[H]
  \centering
  \includegraphics[width=\textwidth]{figures/permApprox/sim_GPD_unconstr_abs_error_shape_by_n-2.pdf}
  \caption{Absolute estimation errors of the GPD shape parameter ($|\hat{\xi} - \xi_0|$) for $\xi_0 = -0.4$ across sample sizes. 
  Boxplots summarize results across 500 replicates per setting, with colors indicating the estimation methods. The starred method (LME*) corresponds to the default estimator in \texttt{permApprox}.
  This figure is adapted from \citep{peschel2026permapprox}.}
  \label{fig:6perm:sim_unconstrained_error_shape}
\end{figure}

As expected, estimation accuracy improves with increasing sample size for all methods, reflecting their asymptotic consistency. 
Across the considered range of shape parameters, the Zhang-Stephens estimator (ZSE) achieved the lowest or near-lowest RMSE in most settings, occasionally matched by the method of moments (MOM) for shape values close to zero. 
The likelihood moment estimator (LME) exhibited comparably high accuracy and closely followed the performance of ZSE across all sample sizes and shape regimes.

Both maximum-likelihood estimators (MLE1D and MLE2D) approached the performance of LME for larger sample sizes ($n \geq 250$) but showed reduced accuracy in small-sample settings ($n \leq 50$). 
The method of moments performed well when $\xi_0 \approx 0$ but deteriorated for larger absolute shape values, particularly in the heavy-tailed regime. 
In contrast, the NLS2 estimator consistently exhibited the largest RMSE values for both parameters.

Overall, ZSE and LME emerged as the most reliable estimators for unconstrained GPD fitting, providing accurate and stable parameter estimates across a wide range of sample sizes and tail behaviors.

%====================================
\subsection{Tolerance calibration for GPD estimation methods}
\label{sec:6perm:param_est:tol}
%====================================

All optimization-based GPD estimators considered in this work rely on a numerical tolerance parameter, denoted by $\tau$ (corresponding to \texttt{tol} in the implementation), that controls the stopping criterion of the underlying estimation procedure. While this parameter is commonly interpreted as a convergence threshold, its exact role differs substantially between methods. 

\textbf{MLE2D:} The parameters are obtained via numerical optimization using the R function \texttt{optim()}. In the unconstrained case, the Nelder-Mead algorithm is employed, where $\tau$ controls convergence based on the relative change in the objective function. More precisely, the optimization terminates when
\begin{equation}
\frac{|f^{(t)} - f^{(t-1)}|}{|f^{(t)}| + \varepsilon} < \tau,
\end{equation}
where $f^{(t)}$ denotes the objective function value at iteration $t$ and $\varepsilon$ is a small constant for numerical stability. In the constrained case, box constraints are imposed and the L-BFGS-B algorithm is used instead. Here, $\tau$ controls convergence via a stopping rule based on the relative reduction of the objective function scaled by machine precision, i.e., the algorithm terminates when further improvements fall below a threshold proportional to $\tau \cdot \varepsilon_{\text{machine}}$.

\textbf{MLE1D:} The one-dimensional MLE is based on a univariate optimization over the endpoint parameter. In our implementation, this optimization is also performed using \texttt{optim()}, and the tolerance therefore plays an analogous role as in MLE2D, controlling the precision with which the profile likelihood is maximized.

\textbf{LME:} The parameter estimation reduces to solving a nonlinear estimating equation for the reparametrized parameter $b$ via a Newton-type iterative scheme. The tolerance governs the stopping criterion of this iterative procedure, typically based on the relative change in successive updates. Hence, $\tau$ determines how accurately the estimating equation is solved.

\textbf{NLS2:} Parameter estimation is again carried out via numerical optimization using \texttt{optim()}, both in the initial transformed problem and in the subsequent refinement step. Consequently, the tolerance controls the convergence of these optimization steps in the same way as for MLE2D, affecting how precisely the least squares criterion is minimized.

\textbf{WNLLSM:} This method also relies on one-dimensional numerical optimization (Brent’s method) for the parameter $b$. The tolerance determines the precision of this optimization and thus the accuracy with which the weighted least squares criterion is minimized.

\textbf{ZSE:} No iterative optimization is performed. Instead, the parameter $b$ is estimated as a weighted average over a predefined grid of candidate values. As a result, this method does not involve a tolerance parameter.

While the tolerance parameter enters the different estimation methods through distinct numerical mechanisms (optimization, root-finding, or grid approximation), it consistently controls the numerical precision of the parameter estimates. In the constrained setting, this precision directly determines how closely the fitted GPD boundary approaches the imposed constraint, and thus influences the fitted density at the evaluation point and the resulting $p$-value approximation.

Consequently, differences in tolerance settings can lead to substantial variation in the estimated tail probabilities across methods, even when the underlying estimation principles are otherwise comparable. This poses a challenge for a fair comparison of estimation methods, as observed differences may be driven by numerical rather than methodological effects.

To address this issue, we conducted a dedicated calibration study to determine suitable tolerance values for each method. The objective of this study is not to optimize parameter recovery, but to harmonize the numerical behavior of the estimators in the constrained regime and thereby ensure comparability in subsequent evaluations.

Specifically, we generated samples from a GPD with shape parameter $\xi = -0.2$ and scale parameter $\sigma = 1$, corresponding to a true boundary at $s = -\sigma/\xi = 5$. 
Each method was used to fit a GPD under the support constraint at an evaluation point 
\begin{equation}
X_{\mathrm{eval}} = 7 > s,
\end{equation}
and the fitted density $f_{\widehat{\mathrm{GPD}}}(X_{\mathrm{eval}})$was evaluated across a range of tolerance values between $10^{-5}$ and $10^{-15}$.

Figure~\ref{fig:6perm:tol_calibration} illustrates that the resulting density estimates can vary substantially with the tolerance parameter, depending on the method. 
In particular, for LME, MLE1D, and NLS2, decreasing the tolerance leads to progressively smaller density values, reflecting convergence towards parameter configurations that place the support boundary increasingly close to the evaluation point. 
In contrast, methods such as WNLLSM exhibit a plateau behavior, where further reduction of the tolerance does not lead to additional changes in the fitted density. 
These differences highlight that the tolerance parameter interacts with the constrained estimation problem in a method-specific manner.

To select appropriate tolerance values, we first identified a numerically stable reference estimator. MLE1D exhibited the smallest variation in $f_{\widehat{\mathrm{GPD}}}(X_{\mathrm{eval}})$ across tolerance values and was therefore used as a reference. For all other methods, we computed the mean absolute error (MAE) of the log-transformed densities at $X_{\mathrm{eval}}$ relative to this reference across simulation repetitions and tolerance values.

Rather than selecting the tolerance value that minimizes the MAE, we applied a plateau criterion: among all tolerance values yielding an MAE within 0.005 of the minimum (on the log scale), we selected the largest tolerance value. 
This strategy favors numerically stable yet computationally efficient configurations, ensuring that further tightening of the tolerance does not lead to meaningful changes in the fitted density.

The selected tolerance values for each method are summarized in Table~\ref{tab:6perm:tol_settings} and were used consistently in all subsequent simulation studies and applications.

\begin{figure}[H]
\centering
\includegraphics[width=\linewidth]{figures/permApprox/tol_calibration_dgpd_vs_rep_all_methods-1.pdf}
\caption{
\textbf{Influence of the tolerance parameter on constrained GPD fitting.}
Fitted densities at the evaluation point ($X_{\mathrm{eval}} = 7$) above the true boundary ($-\sigma/\xi=5$) across simulation repetitions and tolerance values ($10^{-6}$–$10^{-13}$). Each point represents one repetition, colored by tolerance and shaped by estimation method. Highlighted points indicate the selected tolerance values for each method.
This figure is adapted from \citep{peschel2026permapprox}.
}
\label{fig:6perm:tol_calibration}
\end{figure}

\begin{table}[H]
\centering
\caption{Selected tolerance values for the optimization-based GPD parameter estimation methods. The ZSE method is not optimization-based and therefore does not require a tolerance parameter.}
\label{tab:6perm:tol_settings}
\begin{tabular}{ll}
\toprule
\textbf{Method} & \textbf{Selected tolerance} \\
\midrule
LME      & $10^{-7}$  \\
MLE2D    & $10^{-10}$ \\
MLE1D    & $10^{-10}$ \\
NLS2     & $10^{-9}$  \\
WNLLSM   & $10^{-9}$  \\
ZSE      & --- (no optimization) \\
\bottomrule
\end{tabular}
\end{table}


%====================================
\subsection{Simulation study on robustness under support constraints}
\label{sec:6perm:param_est:sim_constrained}
%====================================

We next evaluate the robustness of the estimation methods under the support constraint introduced in Section~\ref{sec:6perm:gpd_tail_modeling}. In contrast to the unconstrained setting, the focus is no longer on parameter recovery, but on the stability and numerical behavior of the estimators when the fitted distribution is forced to extrapolate beyond the observed data.

To this end, data were generated from a bounded GPD with parameters $\sigma_0 = 1$ and $\xi_0 = -0.2$, implying a true upper endpoint at 
\begin{equation}
s_0 = -\frac{\sigma_0}{\xi_0} = 5.
\end{equation}
For each sample size $n \in \{25, 50, 75, 100, 250, 500, 750, 1000\}$, an evaluation point above the true boundary was defined as
\begin{equation}
x_{\mathrm{eval}} = s_0 + \varepsilon, 
\qquad 
\varepsilon \in \{1, 2, 3\},
\end{equation}
corresponding to $x_{\mathrm{eval}} \in \{6, 7, 8\}$. Each estimator was then fitted under the constraint
\begin{equation}
\hat{s} = -\frac{\hat{\sigma}}{\hat{\xi}} > x_{\mathrm{eval}}.
\end{equation}

In this setting, direct comparison of parameter estimates to their true values is not meaningful, as the imposed support constraint introduces a deliberate model misspecification ($s_0 < x_{\mathrm{eval}} < \hat{s}$). Instead, the relevant quantity is the fitted density at the evaluation point,
\begin{equation}
f_{\widehat{\mathrm{GPD}}}(x_{\mathrm{eval}}),
\end{equation}
which directly determines the extrapolated tail probability and thus the resulting $p$-value approximation. A suitable estimator for the \texttt{permApprox} workflow should therefore produce density estimates that are both stable across samples and sufficiently small to yield accurate approximations of extreme tail probabilities.

Figure~\ref{fig:6perm:sim_constrained_dgpd_eps2} shows the fitted densities for $\varepsilon = 2$ (i.e., $x_{\mathrm{eval}} = 7$), which is representative of the overall behavior observed for all considered values of $\varepsilon$. The grouped boxplots display the distribution of $f_{\widehat{\mathrm{GPD}}}(x_{\mathrm{eval}})$ across sample sizes and estimation methods. Only repetitions for which the constraint was active are included, as these correspond to the practically relevant regime of tail extrapolation.

The results reveal substantial differences in the behavior of the estimators under the support constraint. 
While the Zhang-Stephens estimator (ZSE) exhibits the smallest variability in $f_{\widehat{\mathrm{GPD}}}(x_{\mathrm{eval}})$, it consistently produces very large density values across all settings. 
Such inflated densities translate into overly conservative $p$-value approximations and therefore make ZSE unsuitable as a default method in the \texttt{permApprox} framework.

MLE2D shows pronounced variability in both magnitude and spread across sample sizes, indicating instability of the optimization under the imposed constraint. 
In contrast, the remaining methods (WNLLSM, LME, MLE1D, and NLS2) exhibit comparatively stable behavior. 
Within each sample size, the variability follows a consistent ordering,
$\mathrm{SD}_n(\text{WNLLSM}) < \mathrm{SD}_n(\text{LME}) < \mathrm{SD}_n(\text{MLE1D}) < \mathrm{SD}_n(\text{NLS2})$,
although the differences are generally small. The corresponding standard deviations across all values of $\varepsilon$ are summarized in Figure~\ref{fig:6perm:sim_constrained_dgpd_sd}.

\begin{figure}[H]
  \centering
  \includegraphics[width=\textwidth]{figures/permApprox/sim_GPD_evalpoint_dgpd_active_by_n-4.pdf}
  \caption{
    \textbf{Stability of constrained GPD fits at the evaluation point.}
    Fitted densities $f_{\widehat{\mathrm{GPD}}}(x_{\mathrm{eval}})$ at $x_{\mathrm{eval}} = 7$ ($\varepsilon = 2$) across sample sizes $n$ and estimation methods. 
    Only repetitions with active constraint $\hat{s} > x_{\mathrm{eval}}$ are shown. 
    The y-axis is displayed on a $\log_{10}$ scale, while tick labels correspond to the original density values. The starred method (LME*) denotes the default estimator in \texttt{permApprox}.
    This figure is adapted from \citep{peschel2026permapprox}.
  }
  \label{fig:6perm:sim_constrained_dgpd_eps2}
\end{figure}

\begin{figure}[H]
  \centering
  \includegraphics[width=\textwidth]{figures/permApprox/sim_GPD_evalpoint_heat_sd_dgpd-1.pdf}
  \caption{
  Standard deviation of the fitted density values \(f_{\widehat{\mathrm{GPD}}}(x_{\mathrm{eval}})\) across 500 replicates per setting, shown for $\varepsilon \in \{1, 2, 3\}\;\;(\Rightarrow\; x_{\mathrm{eval}} \in \{6, 7, 8\})$. 
  Color intensity reflects standard deviation (SD) of the density: $\sqrt(\mathrm{Var}(f_{\widehat{\mathrm{GPD}}}(x_{\mathrm{eval}})))$, with darker shades indicating higher variability. The starred method (LME*) is the default in \texttt{permApprox}.
  This figure is adapted from \citep{peschel2026permapprox}.
  }
  \label{fig:6perm:sim_constrained_dgpd_sd}
\end{figure}

Overall, these results highlight that robustness under support constraints is primarily determined by the stability of the fitted density at the evaluation point rather than by classical parameter recovery criteria. This perspective is crucial for selecting an appropriate default estimator in the context of permutation $p$-value approximation.

%====================================
\subsection{Runtime comparison of GPD fitting methods}
\label{sec:6perm:param_est:runtime}
%====================================

In addition to statistical accuracy and robustness under support constraints, computational efficiency is an important criterion for selecting a default estimation method, particularly in the context of large-scale permutation testing where GPD fitting is performed repeatedly.

To assess the runtime of the considered methods, we benchmarked all estimators for sample sizes $n \in \{200, 1000\}$ and for both unconstrained and constrained settings. Each configuration was evaluated 50 times per method, and runtimes were recorded in milliseconds (Figure~\ref{fig:6perm:runtime_gpd_methods}).

Overall, the runtime of all methods increases with sample size, as expected. Substantial differences are observed between methods. The NLS2 and WNLLSM methods are by far the slowest approaches across all scenarios. This is consistent with their iterative structure and, in the case of NLS2, the use of a two-step optimization procedure.

The remaining methods exhibit comparable runtime, with the fastest method depending on the specific setting. For smaller sample sizes, LME is the fastest method in the unconstrained case, while ZSE performs best in the constrained setting due to its grid-based implementation. For larger sample sizes, MLE2D is generally the most efficient approach.

Taken together, these results indicate that substantial differences in computational cost exist between estimators, particularly for methods based on nonlinear least squares. This aspect is especially relevant for the \texttt{permApprox} workflow, where GPD fitting is applied repeatedly across many tests. Consequently, methods with high computational cost may become prohibitive in large-scale applications, even if their statistical performance is competitive.

\begin{figure}[H]
  \centering
  \includegraphics[width=\linewidth]{figures/permApprox/gpd_runtime_runtime_gpd_methods-1.pdf}
  \caption{
  \textbf{Runtime comparison of GPD estimation methods.}
  Runtime distributions across sample sizes ($n = 200, 1000$) and boundary settings (with vs.\ without support constraint). Each box represents 50 runs per method and scenario. The starred method (LME*) denotes the default estimator in \texttt{permApprox}.
  This figure is adapted from \citep{peschel2026permapprox}.
  }
  \label{fig:6perm:runtime_gpd_methods}
\end{figure}

%====================================
\subsection{Decision for a default estimation method}
\label{sec:6perm:param_est:decision}
%====================================

In principle, all evaluated estimators are compatible with the \texttt{permApprox} framework. 
However, the choice of a default method requires a careful balance between estimation accuracy, robustness under support constraints, and computational efficiency. The results of the preceding simulation studies provide a comprehensive basis for this decision.

In the unconstrained setting (Section~\ref{sec:6perm:param_est:sim_unconstrained}), ZSE achieved the lowest root mean squared error overall, with LME showing comparably high accuracy across all sample sizes and shape parameters. Thus, both methods can be considered highly competitive in terms of parameter recovery.

In contrast, the constrained setting (Section~\ref{sec:6perm:param_est:sim_constrained}) reveals substantial differences between the estimators. While ZSE exhibits very low variability in the fitted density at the evaluation point, it consistently produces inflated density values, which would translate into overly conservative $p$-value approximations. As accurate estimation of extreme tail probabilities is the primary objective of the \texttt{permApprox} framework, this behavior makes ZSE unsuitable as a default choice.

Among the remaining methods, MLE2D shows considerable variability and unstable behavior under the support constraint. WNLLSM achieves slightly lower variability than LME, but at the cost of substantially increased computational time, being between 15 and 80 times slower in the considered runtime benchmarks. MLE1D performs similarly to LME but tends to exhibit slightly higher variability. NLS2 is consistently the least accurate in the unconstrained setting and among the slowest methods overall.

Taken together, LME provides the most favorable compromise between accuracy, robustness, and computational efficiency. In particular, it combines near-optimal parameter recovery in the unconstrained setting with stable and well-behaved density estimates under the support constraint. It is therefore selected as the default estimator in the \texttt{permApprox} workflow, while the remaining methods are retained as alternative options for comparative analyses or specialized applications.

Finally, we emphasize that all estimation methods are used with the tolerance values determined in the calibration study (Section~\ref{sec:6perm:param_est:tol}). These method-specific tolerance settings are adopted as default throughout the \texttt{permApprox} workflow and are used consistently in all subsequent simulation studies and applications. 
This ensures that differences between estimation methods reflect their intrinsic statistical properties rather than artifacts of numerical convergence.

%%%%%%%%%%%%%%%%%%%%%%%%%%%%%%%%%%%%%
\section{Simulation studies}
\label{sec:6perm:simulations}
%%%%%%%%%%%%%%%%%%%%%%%%%%%%%%%%%%%%%

The previous sections developed the methodological components of the \texttt{permApprox} framework, including threshold selection, the data-adaptive $\varepsilon$ rule, and constrained GPD parameter estimation. We now evaluate the resulting workflow in controlled simulation settings and compare it to existing approaches for approximating permutation $p$-values.

The primary objective of these simulations is to assess the accuracy and stability of $p$-value approximations under finite permutation budgets. Particular emphasis is placed on scenarios in which the true $p$-values are very small, while the number of permutations is limited, resulting in a coarse empirical resolution. Such settings are representative for many applications in microbiome research and beyond, where computational constraints restrict the number of feasible permutations.
 
We consider both a two-sample \textit{t}-test on Gaussian data and a Wilcoxon rank-sum test on exponential data, thereby covering parametric and nonparametric test settings.

%=========================================
\subsection{Compared methods}
\label{sec:6perm:sim:methods}
%=========================================

The compared $p$-value approximation methods are listed below.
For all approaches, tail modeling is applied only when the empirical $p$-value is below $0.1$, ensuring that approximation is restricted to the extreme region of the permutation distribution. For GPD-based methods, empirical $p$-values are used as a fallback whenever the goodness-of-fit test rejects all candidate thresholds.

\begin{itemize}
    \item \textbf{Gamma:} 
    A Gamma distribution is fitted to the full permutation distribution (excluding $T_{\mathrm{obs}}$), following \citet{winkler2016faster}. Parameters are estimated via maximum likelihood, and the $p$-value is obtained as the upper-tail probability of the fitted Gamma distribution evaluated at the observed test statistic..

    \item \textbf{GPD (Knijnenburg):} 
    The classical GPD tail approximation of \citet{knijnenburg2009fewer}. Since the original specification does not clearly define whether one- or two-parameter maximum likelihood estimation is used, we use the one-dimensional version due to its superior numerical stability. For threshold detection, the failure to reject (FTR) procedure with Anderson–Darling (AD) test is used, starting from the top 250 exceedances.
    
    \item \textbf{GPD (Knijnenburg, pow3):} 
    A variant of the above approach in which all test statistics are raised to the third power prior to fitting, as suggested by \citet{knijnenburg2009fewer} to mitigate issues arising from bounded support and negative shape parameters.

    \item \textbf{GPD (Winkler, +obs):}
    \citet{winkler2016faster} compare GPD tail fits both \emph{with} and \emph{without} including the observed test statistic in the permutation distribution. We only evaluate the ``+obs'' variant here as candidate to strictly avoid zero $p$-values. The authors use MOM for parameter estimation, which performs very poorly in some scenarios in our simulations (Section~\ref{sec:6perm:param_est:sim_unconstrained}). We use the standard one-parameter maximum likelihood estimate (MLE1D) instead. For threshold detection, they also use FTR with AD test, but start with the upper quartile of the permutation distribution.

    \item \textbf{GPD (permApprox):} 
    The proposed constrained GPD tail approximation with default configuration as described in Section~\ref{sec:6perm:workflow:default_config}.
\end{itemize}

%=========================================
\subsection{Two-sample \textit{t}-test with Gaussian data}
\label{sec:6perm:sim:ttest}
% Copied from the paper
%========================================= 

To evaluate the accuracy and stability of permutation-based $p$-value approximations under controlled conditions, we first consider a two-sample \textit{t}-test applied to Gaussian data. 
This scenario serves as a benchmark, since the exact theoretical $p$-values are known and Student's \textit{t}-test is optimally suited to the underlying data-generating model. 
It therefore allows a direct assessment of how well each approximation method reproduces the true $p$-values when only a limited number of permutations is available.

We generated two groups of normally distributed data with equal sample sizes $n \in \{25, 50, 100, 250, 500, 1000\}$ and standard deviation $\sigma = 1$. 
For each $n$, we considered effect sizes $d \in \{0.5, 1, 1.5, 2\}$, corresponding to the standardized mean difference between the two groups. 
For each configuration, $B \in \{500, 1000, 5000, 10000\}$ random label permutations were performed to generate the null distribution of the test statistic. 
A total of 1000 independent repetitions were conducted per setting to assess variability. All approximation methods described in Section~\ref{sec:6perm:sim:methods} were applied to the simulated data. 

Figure~\ref{fig:6perm:ttest_ratios_by_n} displays the ratios of approximated permutation $p$-values to the corresponding theoretical \textit{t}-test $p$-values for $d=1$ and $B=1000$ across sample sizes. Ratios close to 1 indicate good agreement with the reference, whereas ratios above 1 correspond to conservative behavior. For small sample sizes, all methods yield ratios centered close to 1. As $n$ increases, however, pronounced differences between the approximation methods emerge.

The Gamma approximation exhibits steadily increasing ratios, reflecting its conservative nature due to modeling the full permutation distribution rather than only the tail. Including the observed statistic in the permutation distribution (``GPD (Winkler, +obs)'') further amplifies this effect, leading to even larger ratios at moderate and large sample sizes. The classical GPD approach of \citet{knijnenburg2009fewer} produces smaller, less conservative $p$-values, but quickly generates zeros even for moderate sample sizes. Already at $n=50$, nearly half of the approximated $p$-values are zero, and for large $n$, almost all vanish to zero due to negative shape estimates (the estimated shape values are provided in Supplementary Figure~\ref{fig:app:6perm:ttest_shapes_by_n}).  Raising all test statistics to the third power, as suggested by the authors, largely prevents zero $p$-values but leads, again, to overly conservative approximations. 	   
In contrast, the proposed \texttt{permApprox} workflow strictly avoids zero $p$-values by construction. While its ratios increase moderately with sample size, they remain substantially closer to 1 than those of the competing methods across all settings, indicating comparatively stable behavior relative to the \textit{t}-test reference.

\begin{figure}[ht]
  \centering
  \includegraphics[width=\textwidth]{figures/permApprox/sim_ttest_papprox_ratios_by_n-1.pdf}
  \caption{\textbf{Two-sample t-test with Gaussian data:} Ratios of approximated to ground-truth $p$-values: $p_{\text{method}} / p_{t\text{-test}}$, stratified by sample size $n$. The horizontal dashed line at 1 indicates perfect agreement with the \textit{t}-test reference. Effect size and number of permutations are fixed to $d=1$ and $B=1000$, respectively. Points indicate individual replicates (1000 per setting). Counts of exact zeros are shown below each group as ``\texttt{0s}=x'' in the corresponding color. For plotting only, zero $p$-values are mapped to a small constant floor so they appear at ``0'' on the $y$-axis. The y-axis is on a log$_{10}$ scale, while tick labels show the original ratios. This figure is adapted from \citep{peschel2026permapprox}}
  \label{fig:6perm:ttest_ratios_by_n}
\end{figure}

The corresponding $p$-values for this setting are shown in Supplementary Figure~\ref{fig:app:6perm:ttest_pvals_by_n}. Additional Supplementary Figures~\ref{fig:app:6perm:ttest_pvals_by_d_and_B} and \ref{fig:app:6perm:ttest_ratios_by_d_and_B} provide $p$-values and ratios stratified by effect size $d$ and number of permutations $B$. The results are consistent across different effect sizes, but slightly change with the number of permutations: With more permutations, the number of zero $p$-values decreases for the classical Knijnenburg approach, but increases when test statistics are raised to the third power. Including the observed test statistic in the null distribution becomes less conservative with increasing $B$. The Gamma approximation is almost unaffected by~$B$.

%=========================================
\subsection{Wilcoxon rank-sum test with exponential data}
\label{sec:6perm:sim:wilcox}
% Copied from the paper
%=========================================

To assess the robustness of the proposed $p$-value approximation method beyond the Gaussian setting, we repeated the simulation study of Section~\ref{sec:6perm:sim:ttest} with the nonparametric Wilcoxon rank-sum test (using the Mann-Whitney $U$ statistic). This setting departs from the assumptions of the \textit{t}-test and involves asymmetric, heavy-tailed data. In this case, inference relies on the permutation null distribution of the Wilcoxon rank-sum statistic, which does not admit a simple parametric form.
We generated two-sample data from exponential distributions with means $1$ and $1+d$, where $d \in \{0, 0.5, 1, 1.5, 2\}$ controls the effect size. 
All five approximation methods considered in the Student's \textit{t}-test setting are applied again here, namely the Gamma approximation, the classical GPD method of \citet{knijnenburg2009fewer}, its power-transformed variant (``pow3''), the ``\,+obs'' variant following \citet{winkler2016faster}, and our proposed \texttt{permApprox} framework.

Figure~\ref{fig:6perm:sim_wilcox_two_panel} shows the results of two different simulation studies.
Panel~(a) shows an illustrative simulation scenario with 1000 tests (200 per effect size), generated for a single fixed data configuration using $n=250$ samples per group and $B=1000$ permutations.
Rather than aggregating results over repeated Monte Carlo replicates, this panel is intended to provide a direct, point-wise comparison of the five approximation methods against the Wilcoxon reference $p$-values across a range of effect sizes.
As in the Gaussian setting, the classical GPD approach produces numerous zero $p$-values for stronger effects due to GPD support violations.
The power-transformed variant largely avoids zero estimates but exhibits pronounced conservativeness.
Both the ``Winkler,\,+obs'' approach and the Gamma approximation avoid zeros by construction, but yield even more conservative $p$-values.
In contrast, \texttt{permApprox} produces smooth approximations that closely track the Wilcoxon $p$-values across all effect sizes, without introducing zeros.

\begin{figure}[H]
  \centering
  \includegraphics[width=\textwidth]{figures/permApprox/sim_wilcox_two_panel.pdf}
  \caption{\textbf{Wilcoxon rank-sum test with exponential data:} Comparison of $p$-value approximations to the Wilcoxon reference.
  \textbf{(a)} Illustrative simulation scenario with 1000 tests across five effect sizes ($d \in \{0, 0.5, 1, 1.5, 2\}$; 200 tests per effect size) and $B=1000$ permutations at $n=250$ samples per group. The x-axis shows the observed Mann-Whitney $U$ statistic, and points represent individual tests. The y-axis displays $p$-values on a log$_{10}$ scale, while tick labels show the original $p$-values.
  \textbf{(b)} Simulation study with 1000 independent replicates per sample size (1 test per replicate) for $d=1$, $B=1000$. The x-axis indicates the per-group sample size, and the y-axis shows ratios of approximated to Wilcoxon $p$-values ($p_{\text{method}} / p_{\text{Wilcoxon}}$) on a log$_{10}$ scale (tick labels show the original ratios). The horizontal dashed line at 1 indicates agreement with the Wilcoxon reference. Points represent individual replicates (1000 per setting).
  In panel~(b), counts of exact zero $p$-values produced by each method are shown below each group as ``\texttt{0s}=x'' in the corresponding color. For visualization only, zero $p$-values are mapped to a small positive constant so they appear at the lower plotting boundary.
  This figure is adapted from \citep{peschel2026permapprox}
  }
  \label{fig:6perm:sim_wilcox_two_panel}
\end{figure}

Panel~(b) presents a comprehensive simulation study with 1000 independent replicates per setting. In each replicate, a single test was conducted, with effect size $d=1$, $B=1000$ permutations, and varying per-group sample sizes from $n=25$ to $n=1000$. Each box summarizes the distribution of ratios of the approximated permutation $p$-values to the corresponding Wilcoxon $p$-values across replicates. 
Ratios close to 1 indicate good agreement with the Wilcoxon reference, while ratios above 1 reflect conservative behavior.
Consistent with the \textit{t}-test results, the classical GPD method produces an increasing number of zero $p$-value estimates as $n$ increases, leading to ratios that collapse to zero for a growing fraction of replicates. 
The remaining approximation methods yield ratios that systematically exceed 1, indicating conservative behavior. 
In contrast, \texttt{permApprox} yields ratios that remain comparatively close to 1 across sample sizes, while strictly avoiding zero $p$-values, suggesting stable behavior across varying test statistic scales and data distributions.
Complementary results, including the corresponding $p$-values for this setting, as well as $p$-values and ratios stratified by $d$ and $B$, are provided in the Supplementary Figures~\ref{fig:app:6perm:wilcox_pvals_by_n} to \ref{fig:app:6perm:wilcox_ratios_by_d_and_B}.


%%%%%%%%%%%%%%%%%%%%%%%%%%%%%%%%%%%%%
\section{Discussion}
\label{sec:6perm:discussion}
%%%%%%%%%%%%%%%%%%%%%%%%%%%%%%%%%%%%%

In this Chapter, we introduced the \texttt{permApprox} framework for accurate permutation $p$-value approximation under limited permutation budgets. Accurate inference in permutation-based testing is often limited by the coarse resolution of empirical permutation $p$-values, which lie on a discrete grid determined by the number of permutations $B$. As a result, extremely small $p$-values cannot be resolved unless $B$ is very large, which is often computationally infeasible. Parametric tail modeling based on the Generalized Pareto Distribution (GPD) has been proposed to overcome this limitation by extrapolating beyond the empirical permutation range \citep{knijnenburg2009fewer}. However, existing GPD-based refinements can yield zero $p$-values when the estimated shape parameter is negative and the fitted tail has a finite upper bound, a problem that has remained largely unresolved. 

The central methodological contribution of this Chapter is a support-constrained GPD tail-fitting approach that ensures that the fitted GPD support extends beyond the observed statistic and thereby avoids structural zero tail probabilities. Specifically, the GPD is fitted under a support constraint, enforced through a data-adaptive $\varepsilon$ rule that balances the trade-off between unrealistically small $p$-values and overly conservative extrapolation. This approach is embedded in a comprehensive workflow that encompasses all essential components of permutation $p$-value approximation, including the choice of the number of initial exceedances, threshold selection, and GPD parameter estimation.

For each component of the workflow, we conducted dedicated simulation studies to identify robust default choices and systematically evaluate competing methods. For the GPD estimation methods, we aimed to balance accuracy in the standard (unconstrained) setting with robustness and numerical stability in the constrained setting. Threshold selection methods were evaluated both with respect to their ability to recover given thresholds and their performance in $p$-value recovery. The resulting default configuration reflects a balance between statistical accuracy, numerical stability, and computational efficiency, tailored to the challenges of permutation testing in high-dimensional settings.

The performance of the combined \texttt{permApprox} workflow was evaluated in two controlled simulation settings reflecting common parametric and nonparametric testing scenarios. In the two-sample \textit{t}-test with Gaussian data, where the theoretical null distribution is known, \texttt{permApprox} closely matched the true $p$-values across a wide range of sample sizes and effect sizes. By contrast, existing GPD-based refinements frequently produced zero $p$-values due to negative shape estimates or became increasingly conservative as sample size increased.

In the nonparametric Wilcoxon rank-sum test with exponential data, which induces highly skewed permutation distributions, constrained GPD fitting again yielded accurate and stable tail probabilities. Competing methods were unable to track the true $p$-values even for moderate test statistics, either collapsing to zero or again being overly conservative. In contrast, \texttt{permApprox} remained robust across all configurations considered in this study, producing smooth and strictly positive approximations even for extreme test statistics. These findings highlight that the benefits of the proposed constrained tail modeling approach are not limited to idealized Gaussian settings but extend to distribution-free tests commonly used in practice.

While no real-data application was presented in this Chapter, the practical relevance of the proposed framework is demonstrated in Chapter~\ref{ch:application} using data from the PASTURE cohort. In this application, \texttt{permApprox} is employed both for differential abundance analysis and for differential network analysis. In particular, the latter setting illustrates a key motivation for the proposed method: network construction procedures are computationally expensive, which typically limits the number of feasible permutations to relatively small values (e.g., $B=1000$). At the same time, differential network analysis involves a large number of simultaneous hypothesis tests. By improving the resolution of small $p$-values under these constraints, \texttt{permApprox} increases the power to detect differential associations after adjusting for multiplicity and enables more reliable inference in such high-dimensional settings.
Despite its advantages, \texttt{permApprox} has a few practical limitations. First, the approach assumes that the upper tail of the permutation distribution can be adequately approximated by a GPD and that a suitable threshold can be identified. In settings where the permutation tail is highly irregular or contains too few exceedances, the goodness-of-fit test may reject the GPD model, in which case tail approximation is inappropriate. In such cases, empirical permutation $p$-values remain the appropriate fallback, and alternative approaches such as modeling the full permutation distribution may be preferable.

Second, although the proposed data-adaptive $\varepsilon$ rule performed robustly across both Gaussian and nonparametric simulation settings, it necessarily influences the balance between overly optimistic and overly conservative tail probabilities. In rare cases with exceptionally extreme test statistics, the constraint may require additional refinement to prevent numerical underflow, at the cost of slightly more conservative $p$-values for other tests.

Finally, constrained tail fitting introduces some computational overhead compared to relying solely on empirical permutation $p$-values. However, this additional cost is modest in typical analysis settings and negligible compared to the computational burden of increasing the number of permutations.

Taken together, the results of this Chapter show that \texttt{permApprox} successfully eliminates the long-standing problem of zero $p$-values in GPD-based permutation $p$-value approximation and provides a broadly applicable workflow for refining permutation tails. Importantly, the method is not tied to any specific test statistic or application domain. Any permutation-based procedure that returns (i) a vector of observed test statistics and (ii) the corresponding matrix of permuted statistics can be directly extended with the proposed workflow. By improving the reliability and resolution of small $p$-values, \texttt{permApprox} broadens the scope of accurate and computationally feasible permutation-based inference across a wide range of data-analytic applications.